\documentclass[aspectratio=169, xcolor=dvipsnames]{beamer}

\AtBeginSection[]
{
    \begin{frame}
        \frametitle{Visão Geral}
        \tableofcontents[currentsection]
    \end{frame}
}

\usefonttheme{professionalfonts} % using non standard fonts for beamer
\usefonttheme{serif} % default family is serif
\usepackage{fontspec}
\setmainfont{XCharter}[
    UprightFont    = *-Roman,
    BoldFont       = *-Bold,
    ItalicFont     = *-Italic,
    BoldItalicFont = *-BoldItalic,
]

\usepackage{hyperref}
\usepackage{graphicx} % Allows including images
\usepackage{booktabs} % Allows the use of \toprule, \midrule and \bottomrule in tables
\input{./slides/SimplePlusTheme/beamerthemeSimplePlus.sty}
\input{./slides/SimplePlusTheme/beamerinnerthemeSimplePlus.sty}
\input{./slides/SimplePlusTheme/beamerfontthemeSimplePlus.sty}
\input{./slides/SimplePlusTheme/beamercolorthemeSimplePlus.sty}

\usepackage{listings}
\setbeamercovered{transparent}

\usepackage{biblatex}
\addbibresource{slides/reference.bib}
\renewcommand*{\bibfont}{\footnotesize}
% Tamanho dos slides de "Leituras Recomendadas": maior que a bibliografia
% completa (\bibfont acima, que rende \tiny via a redefinição de \footnotesize)
% para destacar as leituras curadas.
\newcommand{\leiturasize}{\small}

\usepackage[brazilian ]{babel}
\usepackage{csquotes}
\usepackage{pgfplots}
\pgfplotsset{compat=1.18}

\usepackage{emoji}
\usepackage{amsmath}
\usepackage[xcharter,bigdelims,vvarbb]{newtxmath}
% newtxmath's operators font requests the fontspec family `XCharter(0)' in OT1;
% without these substitutions math digits and operator names silently fall back
% to Computer Modern (LaTeX Font Warning: Font shape `OT1/XCharter(0)/m/n' undefined).
\DeclareFontFamily{OT1}{XCharter(0)}{}
\DeclareFontShape{OT1}{XCharter(0)}{m}{n}{<->ssub * XCharter-TLF/m/n}{}
\DeclareFontShape{OT1}{XCharter(0)}{m}{it}{<->ssub * XCharter-TLF/m/it}{}
\DeclareFontShape{OT1}{XCharter(0)}{b}{n}{<->ssub * XCharter-TLF/b/n}{}
\DeclareFontShape{OT1}{XCharter(0)}{bx}{n}{<->ssub * XCharter-TLF/b/n}{}

\usepackage{bm}
\usepackage[table]{xcolor}
\usepackage{xcolor}
\usepackage{algpseudocode}

\usepackage{cancel}
\usepackage{ulem}

\usetikzlibrary{shapes}
\usetikzlibrary{arrows}
\usetikzlibrary {arrows.meta}
\usetikzlibrary{calc,positioning}
\usepgfplotslibrary{fillbetween}
\usetikzlibrary{angles,quotes}
\usetikzlibrary{tikzmark}


\definecolor{PastelRed}{HTML}{FFC1CC}
\definecolor{PastelBlue}{HTML}{C2F0FF}
\definecolor{PastelBlue2}{HTML}{3182BD}

\definecolor{PastelPink}{HTML}{FFCBE1}
\definecolor{PastelGreen}{HTML}{D6E5BD}
\definecolor{PastelYellow}{HTML}{F9E1A8}
\definecolor{PastelBlue}{HTML}{BCD8EC}
\definecolor{PastelPurple}{HTML}{DCCCEC}
\definecolor{PastelOrange}{HTML}{FFDAB4}


\renewcommand{\footnotesize}{\tiny}

\newcommand{\mespor}{
  \ifcase\month
    \or janeiro
    \or fevereiro
    \or março
    \or abril
    \or maio
    \or junho
    \or julho
    \or agosto
    \or setembro
    \or outubro
    \or novembro
    \or dezembro
  \fi
}
% \usepackage{csquotes}
% \usepackage{minted}
% \usemintedstyle{xcode}
\definecolor{vscLightBackground}{HTML}{FFFFFF}
\definecolor{vscLightText}{HTML}{000000}
\definecolor{vscLightGreen}{HTML}{008000}        % For Comments AND Numbers
\definecolor{vscLightRed}{HTML}{A31515}          % For Strings
\definecolor{vscLightTeal}{HTML}{267F99}
\definecolor{vscBlue}{HTML}{0000FF}% For main keywords (if, for, def...)
\definecolor{vscLightPytorch}{HTML}{800080}      % For PyTorch keywords (dark magenta)

\lstdefinestyle{vscode-light}{
    language=Python,
    backgroundcolor=\color{vscLightBackground},
    basicstyle=\ttfamily\scriptsize\color{vscLightText},
    keywordstyle=\color{vscLightTeal},
    commentstyle=\color{vscLightGreen},
    stringstyle=\color{vscLightRed},
    % Style for PyTorch keywords (using keyword group 2)
    keywordstyle=[2]\color{vscLightPytorch}, 
    morekeywords={self, True, False, None},
    morekeywords=[2]{
        torch, Tensor, device, dtype, to, nn, optim, autograd, utils, data,
        Module, Linear, Conv2d, ReLU, CrossEntropyLoss, Adam, SGD,
        Dataset, DataLoader, randn, zeros, ones, arange, cat, stack,
        view, reshape, sum, mean, max, min, matmul, no_grad, backward
    },
    keywordstyle=[3]\color{vscBlue}, 
    morekeywords=[3]{
        model, loss, y, y_hat, x , optimizer
    },
    frame=single,
    tabsize=1,
    showstringspaces=false,
    breaklines=true,
    captionpos=b,
    extendedchars=true,
    numbers=left,
}

\usepackage[cache=true]{minted}
% minted v3+ (TeX Live 2024+) auto-detects the output directory.
% minted v2 (Ubuntu apt TeX Live 2023) needs it set explicitly to
% match latexmkrc's $out_dir = 'build'.
\makeatletter
\@ifpackagelater{minted}{2024/01/01}{}{%
  \def\minted@outputdir{build/}%
}
\makeatother
\setminted{
    fontsize=\scriptsize,
    style=vs,
    breaklines=true,
    numbers=left,
    numbersep=5pt,
    frame=single,
}
\providecommand{\citeauthoryear}[1]{(\citeauthor{#1},~\citeyear{#1})}

\providecommand{\thetav}{\bm{\theta}}

% vetor coluna 2D com colchetes
\providecommand{\cvec}[2]{\begin{bmatrix} #1 \\ #2 \end{bmatrix}}

\providecommand{\varvector}[1]{\mathbf{#1}}
% instance
\providecommand{\X}{\varvector{X}}
\providecommand{\mlinstance}{\varvector{X}}
\providecommand{\mlinstancei}{\varvector{X}^{(i)}}
\providecommand{\mlinstancex}[1]{\varvector{X}^{(#1)}}

% target
\providecommand{\y}{\varvector{y}}
\providecommand{\mltarget}{\varvector{y}}
\providecommand{\mltargeti}{\varvector{y}^{(i)}}
\providecommand{\mltargetx}[1]{\varvector{y}^{(#1)}}

%hats
\providecommand{\yhat}{\hat{\varvector{y}}}
\providecommand{\mltargethat}{\hat{\varvector{y}}}
\providecommand{\mltargethati}{\hat{\varvector{y}}^{(i)}}
\providecommand{\mltargethatx}[1]{\hat{\varvector{y}}^{(#1)}}

% linear reg
\providecommand{\mllinearregression}{\mltargethat=\mlinstance\bm{\theta}}

\providecommand{\tikzarrowcolor}[3]{%
  \begin{tikzpicture}[remember picture, overlay]
    % Arrow body
    \path[fill=#3] ([shift={(#1,#2)}]current page.south west) 
      rectangle ++(0.3, -0.2);
    % Arrow head
    \path[fill=#3] ([shift={(#1+0.3,#2+0.1)}]current page.south west) -- 
         ++(0, -0.4) -- 
         ++(0.2, 0.2) -- 
         cycle;
  \end{tikzpicture}%
}

\providecommand{\tikzarrow}[2]{%
    \tikzarrowcolor{#1}{#2}{Red}
}



\DeclareMathOperator*{\argmax}{argmax}
\DeclareMathOperator*{\argmin}{argmin}



%----------------------------------------------------------------------------------------
%    EXTRA COLORS / COMMANDS FOR THIS LECTURE
%----------------------------------------------------------------------------------------
\definecolor{LightBlue}{HTML}{DBEAFE}
\definecolor{DarkBlue}{HTML}{1565C0}
\definecolor{DarkRed}{HTML}{C62828}
\definecolor{LightGray}{HTML}{F5F5F5}
\definecolor{MidGray}{HTML}{757575}
\definecolor{CellGreen}{HTML}{C8E6C9}
\definecolor{CellYellow}{HTML}{FFF9C4}
\definecolor{CellRed}{HTML}{FFCDD2}
\definecolor{CellBlue}{HTML}{BBDEFB}
\definecolor{CellPurple}{HTML}{E1BEE7}
\definecolor{DarkGreen}{HTML}{1B5E20}
\definecolor{DarkOrange}{HTML}{E65100}
\definecolor{LightOrange}{HTML}{FFE0B2}

%----------------------------------------------------------------------------------------
%    TITLE PAGE
%----------------------------------------------------------------------------------------

\title{Deep Learning}
\subtitle{\textit{Treinamento de LLMs: do Pré-treinamento ao Alinhamento}}

\author{Lucas Silveira Kupssinskü}

\institute
{ Machine Learning Theory and Applications Laboratory \\ PUCRS
}
\date{\mespor de \the\year}

%----------------------------------------------------------------------------------------
%    PRESENTATION SLIDES
%----------------------------------------------------------------------------------------

\begin{document}
  \begin{frame}
    \titlepage
  \end{frame}

  \begin{frame}{Visão Geral}
    \tableofcontents
  \end{frame}

  %======================================================================================
  % SECTION 1: PRÉ-TREINAMENTO
  %======================================================================================
  \section{Pré-treinamento}

  %--- Slide: Objetivo do pré-treinamento ---
  \begin{frame}{Objetivo: \textit{Next-Token Prediction}}
          O pré-treinamento de LLMs se baseia em \textbf{modelagem de linguagem causal}:
        dado um prefixo de tokens, prever o próximo.
    \begin{columns}
      \begin{column}{0.52\textwidth}

        \begin{block}{\textit{Loss Function}: \textit{Cross Entropy}}
          \[
            \mathcal{L}_{\text{LM}}(\theta) = -\frac{1}{T}\sum_{t=1}^{T}
            \log P_\theta(x_t \mid x_1, \ldots, x_{t-1})
          \]
        \end{block}

      \end{column}
      \begin{column}{0.45\textwidth}


        \centering
        \resizebox{\linewidth}{!}{%
        \begin{tabular}{ccc}
          \toprule
          \textbf{Contexto} & $\to$ & \textbf{Alvo} \\
          \midrule
          \texttt{O gato} & $\to$ & \texttt{sentou} \\
          \texttt{O gato sentou} & $\to$ & \texttt{no} \\
          \texttt{O gato sentou no} & $\to$ & \texttt{tapete} \\
          \bottomrule
        \end{tabular}}


      \end{column}
    \end{columns}
    \vspace{1em}
    
    Uma única sequência de comprimento $T$ gera $T$ pares (contexto, alvo); o primeiro tem contexto vazio (na prática, apenas o token inicial \texttt{<s>}).

    \vspace{1em}
    \textbf{Propriedade chave:} o corpus de texto é a própria supervisão: não há necessidade de anotação humana.
  \note{
    \textbf{Intuição da equação (perda LM):}
    \begin{itemize}
      \item A perda é simplesmente a \textbf{entropia cruzada} média sobre os $T$ tokens da sequência: o modelo deve maximizar a probabilidade do token correto em cada posição.
      \item O fator $-1/T$ normaliza pelo comprimento, permitindo comparar perdas entre sequências de tamanhos diferentes.
      \item \textbf{Teacher forcing:} durante o treino, o contexto fornecido ao modelo é sempre o texto \emph{real} (não as predições anteriores do modelo). Isso permite processar todos os $T$ pares em \emph{paralelo} com uma única passagem pelo transformador (é por isso que o pré-treinamento de LLMs é viável em escala).
      \item Ponto para enfatizar: não há anotadores humanos! O texto em si é tanto a entrada quanto o rótulo. Isso é o que torna viável treinar em trilhões de tokens.
    \end{itemize}
  }
  \end{frame}

  %--- Slide: Datasets ---
  \begin{frame}{Datasets de Pré-treinamento}
        Grandes corpora públicos
          \begin{itemize}
            \item \textbf{Common Crawl}: rastreamento da web, bruto e ruidoso
            \item \textbf{C4} (\textit{Colossal Clean Crawled Corpus}): versão filtrada do CC
            \item \textbf{The Pile}: 825 GB de texto diversificado (livros, código, arXiv, GitHub)
            \item \textbf{ROOTS} (BigScience): corpus multilíngue (46 idiomas, 1.6 TB)
            \item \textbf{FineWeb} (HuggingFace): 15 T de tokens filtrados da web
          \end{itemize}
        Etapas de processamento
          \begin{enumerate}
            \item \textbf{Filtragem}: remoção de conteúdo ruim e spam
              (heurísticas + classificadores)
            \item \textbf{Deduplicação}: remoção de documentos duplicados ou
              quase-duplicados
            \item \textbf{Tokenização}: conversão para tokens
            \item \textbf{Shuffling}: evitar correlações temporais no gradiente
          \end{enumerate}
        \textcolor{DarkRed}{\small \textbf{Regra geral:} qualidade $>$ quantidade: um corpus menor e limpo supera um corpus bruto maior.}
  \note{
    \textbf{Pontos para a figura/tabela de datasets:}
    \begin{itemize}
      \item Contextualizar a escala: The Pile tem 825 GB, FineWeb tem 15 trilhões de tokens (ambos cabem num disco SSD, mas o treino consome meses em centenas de GPUs).
      \item \textbf{Deduplicação} é crítica: documentos duplicados causam ``memorização'' e \textit{overfitting}. MinHash é $O(N)$ por documento; suffix arrays encontram duplicatas exatas em $O(N\log N)$.
      \item \textbf{Shuffling} embaralha a ordem temporal dos documentos: sem isso, o gradiente tem correlação serial (notícias do mesmo dia juntas, por exemplo) e o Adam acumula estimativas enviesadas.
      \item A regra qualidade $>$ quantidade pode surpreender os alunos: intuitivamente mais dados = melhor. Mostre o exemplo do LIMA na próxima seção para reforçar esse ponto.
    \end{itemize}
  }
  \end{frame}

  %--- Slide: Tokenização ---
  \begin{frame}{Recapitulação: Tokenização e Vocabulário}
    A tokenização (vista em detalhe na aula de tokenização) converte texto bruto em sequências de \textit{tokens}. O que importa para o treinamento:

    \vspace{0.4em}
    \begin{block}{Tamanho do vocabulário e impacto}
      \begin{itemize}
        \item Vocabulário menor $\Rightarrow$ sequências mais longas, mais computação por texto
        \item Vocabulário maior $\Rightarrow$ matriz de \textit{embedding} e projeção de saída maiores, cobertura melhor de outras línguas e código
        \item Modelos modernos: 32\,k -- 200\,k tokens (BPE ou Unigram)
      \end{itemize}
    \end{block}

    \vspace{0.4em}
    O \textit{tokenizer} é \textbf{treinado antes} do modelo, no mesmo corpus, e permanece fixo durante o treinamento: trocar ou estender o vocabulário depois exige retreinar o \textit{tokenizer} e ajustar os \textit{embeddings}.
  \note{
    \textbf{Pontos para a recapitulação:}
    \begin{itemize}
      \item Exemplos de tamanho de vocabulário, se perguntarem: GPT-2 usa 50\,257 tokens (256 bytes base + 50\,000 \textit{merges} do BPE + 1 token especial \texttt{<|endoftext|>}); LLaMA-2 usa 32\,000; LLaMA-3 quadruplicou para 128\,256, principalmente para melhorar a cobertura de idiomas não-ingleses e código.
      \item \textbf{Trade-off prático:} vocabulário $2\times$ maior $\Rightarrow$ embedding de entrada e projeção de saída $2\times$ maiores (para um modelo de 7B, isso pode representar 3--5\% do total de parâmetros).
      \item Lembrar que o tokenizador é treinado \emph{antes} e permanece \emph{congelado}: se você quiser adicionar uma língua nova depois, precisa retreinar o tokenizador e fazer \textit{embedding extension}.
    \end{itemize}
  }
  \end{frame}

  %--- Slide: Schedule do LR (warm-up + cosine + WSD) ---
  \begin{frame}{Pré-treinamento: \textit{Schedule} da Taxa de Aprendizado \cite{loshchilov2019adamw}}
    \small
    \begin{columns}[T]
      \begin{column}{0.52\textwidth}
        \begin{block}{1. \textit{Warm-up} linear}
          \begin{itemize}
            \item LR cresce de 0 até $\eta_{\max}$ em $T_{\text{wu}}$ passos (1--2\% do treino)
            \item Os momentos do Adam começam em zero: sem \textit{warm-up}, as primeiras atualizações seriam enormes
          \end{itemize}
        \end{block}
        \begin{block}{2. \textit{Cosine decay}}
          De $\eta_{\max}$ a $\eta_{\min}$ (tipicamente $0.1\,\eta_{\max}$), para $T_{\text{wu}} \le t \le T$:
          \vspace{-0.3em}
          \[
            \eta_t = \eta_{\min} + \tfrac{\eta_{\max}-\eta_{\min}}{2}
              \!\left(1+\cos\tfrac{\pi\,(t - T_{\text{wu}})}{T - T_{\text{wu}}}\right)
          \]
        \end{block}
      \end{column}
      \begin{column}{0.45\textwidth}
        \centering
        \begin{tikzpicture}
          \begin{axis}[
            width=5.8cm, height=3.5cm,
            xlabel={passo $t$},
            ylabel={$\eta_t$},
            xmin=0, xmax=1,
            ymin=0, ymax=1.18,
            xtick={0, 0.1, 1},
            xticklabels={$0$, $T_{\text{wu}}$, $T$},
            ytick={0.1, 1.0},
            yticklabels={$\eta_{\min}$, $\eta_{\max}$},
            font=\scriptsize,
            axis lines=left,
            tick align=outside,
          ]
            % Warm-up phase (linear, blue)
            \addplot[DarkBlue, thick, domain=0:0.1, samples=20] {10*x};
            \addplot[DarkBlue!20, fill, opacity=0.4, domain=0:0.1, samples=20]
              {10*x} \closedcycle;
            % Cosine decay phase (red)
            \addplot[DarkRed, thick, domain=0.1:1.0, samples=80]
              {0.1 + 0.45*(1 + cos(180*(x-0.1)/0.9))};
            \addplot[DarkRed!20, fill, opacity=0.4, domain=0.1:1.0, samples=80]
              {0.1 + 0.45*(1 + cos(180*(x-0.1)/0.9))} \closedcycle;
            % WSD alternative (dashed gray)
            \addplot[MidGray, thick, dashed, domain=0.1:0.85, samples=2] {1.0};
            \addplot[MidGray, thick, dashed, domain=0.85:1.0, samples=2] {1.0 - 6*(x-0.85)};
            % Phase labels
            \node[font=\tiny, text=DarkBlue, anchor=west]
              at (axis cs:0.11, 0.35) {warm-up};
            \node[font=\tiny, text=DarkRed]
              at (axis cs:0.55, 0.72) {cosine decay};
            \node[font=\tiny, text=MidGray]
              at (axis cs:0.50, 1.08) {WSD (tracejado)};
          \end{axis}
        \end{tikzpicture}

        {\footnotesize\raggedright \textbf{Alternativa: WSD} (\textit{Warmup-Stable-Decay}) \cite{hu2024minicpm}: LR constante na maior parte do treino e decaimento curto no fim (10--20\% dos passos). Permite estender o treino sem fixar $T$ de antemão.\par}
      \end{column}
    \end{columns}
  \note{
    \textbf{Warm-up, cosine e WSD:}
    \begin{itemize}
      \item \textbf{Warm-up:} os estimadores de média e variância do Adam são inicializados em zero, portanto as primeiras atualizações seriam muito grandes sem warm-up. O warm-up linear dá tempo para o Adam ``calibrar'' seus estimadores.
      \item \textbf{Cosine:} em $t=T_{\text{wu}}$: $\cos(0) = 1 \Rightarrow \eta = \eta_{\max}$. Em $t=T$: $\cos(\pi) = -1 \Rightarrow \eta = \eta_{\min}$. LLaMA usa $\eta_{\min} = 0.1\,\eta_{\max}$. A transição suave evita as instabilidades que os degraus do step decay podem causar.
      \item \textbf{Limitação do cosine:} a curva depende de $T$. Se decidirmos treinar por mais tokens depois, o schedule inteiro muda. O WSD resolve isso: mantém o LR constante, guarda checkpoints, e faz um decaimento curto a partir de qualquer checkpoint para ``colher'' um modelo. Também facilita experimentos de \textit{scaling law}: o mesmo run produz modelos com vários orçamentos de tokens.
    \end{itemize}
  }
  \end{frame}

  %--- Slide: Otimizador, clipping e batch ---
  \begin{frame}{Pré-treinamento: Otimizador, \textit{Gradient Clipping} e \textit{Batch Size}}
    \small
    \begin{columns}[T]
      \begin{column}{0.50\textwidth}
        \begin{block}{Otimizador padrão: AdamW}
          \begin{itemize}
            \item $\beta_1{=}0.9$, $\beta_2{=}0.95$, $\epsilon{=}10^{-8}$, \textit{weight decay} $\lambda{=}0.1$
            \item $\beta_2{=}0.95$ (e não $0.999$): histórico mais curto, reage mais rápido a mudanças no gradiente
          \end{itemize}
        \end{block}
        \begin{block}{\textit{Gradient clipping} por norma global}
          Seja $\mathbf{g}$ o gradiente de \textbf{todos os parâmetros da rede} concatenados em um único vetor, antes de entrar no Adam.
          Se $\|\mathbf{g}\|_2 > 1.0$, escala-se o vetor inteiro: $\mathbf{g} \leftarrow \mathbf{g}/\|\mathbf{g}\|_2$.
          Preserva a direção (ao contrário do \textit{clipping} por valor) e contém picos de \textit{loss}.
        \end{block}
      \end{column}
      \begin{column}{0.47\textwidth}
        \begin{block}{\textit{Batch ramp-up} \cite{mccandlish2018batch}}
          \begin{itemize}
            \item \textit{Batch} cresce ao longo do treino (ex.: de 1M para 4M tokens)
            \item O \textbf{ruído do gradiente cresce} conforme a \textit{loss} cai, e com ele o \textit{batch} eficiente (\textit{critical batch size})
            \item Início: \textit{batches} pequenos aproveitam melhor cada token
            \item Fim: \textit{batches} grandes ocupam mais GPUs sem desperdiçar amostras
          \end{itemize}
        \end{block}
        \textcolor{DarkRed}{\textit{Batch} grande não melhora a generalização por si só: a motivação é eficiência computacional.}
      \end{column}
    \end{columns}
  \note{
    \textbf{Otimizador, clipping e batch:}
    \begin{itemize}
      \item \textbf{$\beta_2 = 0.95$} (vs.\ padrão 0.999): acumula histórico de gradientes mais curto, tornando o Adam mais adaptável a mudanças rápidas de gradiente comuns no pré-treinamento. Com $0.999$, um pico isolado de gradiente ``contamina'' a estimativa de variância por milhares de passos.
      \item \textbf{Gradient clipping} (norma $\leq 1.0$): a norma é calculada sobre o gradiente de todos os parâmetros do modelo de uma vez (bilhões de entradas), não camada a camada, e o mesmo fator de escala é aplicado a todos. Acontece depois do backward (e do \textit{unscale} do AMP) e antes do \texttt{optimizer.step()}: o que entra nos momentos $m$ e $v$ do Adam já é o gradiente cortado, então um pico isolado não contamina o histórico. Previne explosão de gradiente especialmente no início do treino. Diferenciar de gradient clipping por valor (clip element-wise), que distorce a direção do gradiente. Observar que, como o Adam normaliza por $\sqrt{v}$, o clipping por norma não limita o tamanho do passo, só o que alimenta o otimizador.
      \item \textbf{Critical batch size} (McCandlish et al., 2018): a razão entre a variância do gradiente e o quadrado da sua norma (\textit{gradient noise scale}) define o batch a partir do qual aumentar $B$ deixa de reduzir o número de passos necessários. Essa escala cresce ao longo do treino, por isso o ramp-up. Batch grande demais no início desperdiça amostras (retornos decrescentes); batch pequeno demais no fim desperdiça GPUs.
      \item Cuidado com a intuição ``batch grande generaliza melhor'': a evidência clássica (Keskar et al., 2017) vai na direção oposta em visão. Em LLMs o efeito é pequeno se o LR for ajustado; o argumento é de custo, não de generalização.
    \end{itemize}
  }
  \end{frame}

  %--- Slide: Receitas reais de pré-treinamento ---
  \begin{frame}{Na Prática: Receitas de Pré-treinamento de Modelos Reais}
    \centering
    \resizebox{0.74\linewidth}{!}{%
    \begin{tabular}{lccc}
      \toprule
      & \textbf{LLaMA-7B} \cite{touvron2023llama} & \textbf{Llama 3 405B} \cite{grattafiori2024llama3} & \textbf{SmolLM2-1.7B} \cite{allal2025smollm2} \\
      \midrule
      Tokens de treino     & 1,0T & 15,6T & 11T \\
      \textit{Batch} (tokens) & 4M & 4M $\to$ 8M $\to$ 16M & 2M \\
      Contexto             & 2\,048 & 4\,096 $\to$ 8\,192 & 2\,048 \\
      $\eta_{\max}$        & $3\times10^{-4}$ & $8\times10^{-5}$ & $5\times10^{-4}$ \\
      \textit{Warm-up}     & 2\,000 passos & 8\,000 passos & 2\,000 passos \\
      \textit{Schedule}    & \textit{cosine} & \textit{cosine} & WSD \\
      $\eta_{\min}$        & $0{,}1\,\eta_{\max}$ & $0{,}01\,\eta_{\max}$ & $0$ (linear, últimos 10\%) \\
      Otimizador           & AdamW, $\beta_2{=}0.95$ & AdamW & AdamW, $\beta_2{=}0.95$ \\
      \bottomrule
    \end{tabular}}

    \vspace{1.1em}
    \begin{columns}[T]
      \begin{column}{0.48\textwidth}
        \textcolor{DarkBlue}{\textbf{O que se repete}}
        \begin{itemize}\small
          \item \textit{Warm-up} curto (menos de 1\% dos passos) e decaimento suave
          \item $\beta_2{=}0.95$ e \textit{batches} de milhões de tokens em todos
          \item Contexto curto no pré-treino; contextos longos vêm numa fase final de extensão
        \end{itemize}
      \end{column}
      \begin{column}{0.48\textwidth}
        \textcolor{DarkRed}{\textbf{O que muda com a escala}}
        \begin{itemize}\small
          \item $\eta_{\max}$ cai conforme o modelo cresce: $5\times10^{-4}$ (1,7B) até $8\times10^{-5}$ (405B)
          \item O \textit{batch} cresce durante o treino nos modelos maiores (\textit{ramp-up} do Llama 3)
          \item WSD quando o orçamento de tokens não é fixado de antemão (SmolLM2)
        \end{itemize}
      \end{column}
    \end{columns}
  \note{
    \textbf{Como usar a tabela:}
    \begin{itemize}
      \item \textbf{LLaMA 1} (Touvron et al., 2023, seção 2.3): AdamW com $\beta_1=0.9$, $\beta_2=0.95$, \textit{cosine} até 10\% do LR máximo, \textit{weight decay} 0.1, \textit{gradient clipping} 1.0, 2\,000 passos de \textit{warm-up}. O LR e o batch variam com o tamanho: 7B usa $3\times10^{-4}$ e 65B usa $1{,}5\times10^{-4}$, ambos com 4M tokens por batch. É a origem dos ``valores padrão'' citados no slide anterior de otimizador.
      \item \textbf{Llama 3 405B} (Grattafiori et al., 2024, seção 3.4.1): LR de pico $8\times10^{-5}$, \textit{warm-up} linear de 8\,000 passos, \textit{cosine} até $8\times10^{-7}$ ao longo de 1,2M passos. O batch começa em 4M tokens com sequências de 4\,096, dobra para 8M tokens e sequências de 8\,192 após 252M tokens, e dobra de novo para 16M após 2,87T tokens. A justificativa dos autores é estabilidade no início do treino: é exatamente o \textit{batch ramp-up} do slide anterior. \textit{Weight decay} e \textit{clipping} não são informados no relatório, por isso a célula só diz AdamW.
      \item \textbf{SmolLM2-1.7B} (Allal et al., 2025): 11T tokens em quatro estágios com misturas de dados diferentes, LR de pico $5\times10^{-4}$, 2\,000 passos de \textit{warm-up}, batch de 2M tokens, sequência de 2\,048 (estendida para 8\,192 no fim). O \textit{schedule} é WSD: LR constante e decaimento linear a zero nos últimos 10\% do treino. Como o treino foi planejado em estágios, o WSD permitiu ajustar a mistura de dados sem reiniciar o \textit{schedule}.
      \item Enfatizar que os números não são universais: modelos com arquiteturas ou tokenizadores diferentes (Qwen, Gemma, DeepSeek) usam outras combinações, mas a estrutura (\textit{warm-up}, decaimento, batch grande, $\beta_2$ menor) é a mesma.
    \end{itemize}
  }
  \end{frame}

  %======================================================================================
  % SECTION 2: FINE-TUNING SUPERVISIONADO
  %======================================================================================
  \section{Fine-tuning Supervisionado (SFT)}

  %--- Slide: Objetivo do SFT ---
  \begin{frame}{SFT: Ajuste para Seguir Instruções}
        Após o pré-treinamento, o modelo sabe ``completar texto'', mas não necessariamente \textbf{seguir instruções} de forma útil.
        
        \vspace{0.2em}
        
        O \textit{Supervised Fine-Tuning} (SFT), ou \textit{Instruction Tuning},
        ajusta o modelo em pares \texttt{(instrução, resposta)} de alta qualidade.
    \begin{columns}
      \begin{column}{0.52\textwidth}

        \vspace{0.2em}
        \begin{block}{Função de perda do SFT}
          \[
            \mathcal{L}_{\text{SFT}}(\theta) = -\sum_{t \in \mathcal{R}}
            \log P_\theta(x_t \mid x_{<t})
          \]
          onde $\mathcal{R}$ são \textbf{apenas os tokens da resposta}
          (não do prompt/instrução).
        \end{block}

        \textcolor{DarkRed}{\small É a \textbf{mesma perda do pré-treino}, mas mascarada: o gradiente vem só dos tokens que o modelo deve gerar, não dos que ele recebe do usuário.}
        % Note: \note command placed at end of frame
      \end{column}
      \begin{column}{0.45\textwidth}
        \textcolor{DarkBlue}{\textbf{Exemplo de amostra de SFT}}
        \vspace{0.3em}

        \begin{tikzpicture}
          \node[draw=DarkBlue, rounded corners, fill=LightBlue, text width=5.5cm,
                font=\small, inner sep=6pt] (prompt) {
            \textcolor{MidGray}{\textbf{[INSTRUÇÃO: mascarado]}}\\
            Explique o que é retropropagação.
          };
          \node[draw=DarkGreen, rounded corners, fill=CellGreen, text width=5.5cm,
                font=\small, inner sep=6pt, below=0.2cm of prompt] (resp) {
            \textcolor{DarkGreen}{\textbf{[RESPOSTA: perda calculada aqui]}}\\
            Retropropagação é o algoritmo que calcula gradientes\ldots
          };
        \end{tikzpicture}
      \end{column}
    \end{columns}
  \note{
    \textbf{Intuição da perda SFT e da figura:}
    \begin{itemize}
      \item A perda SFT é matematicamente idêntica à perda de pré-treinamento (NLL / entropia cruzada), mas com uma \textbf{máscara}: o gradiente é zerado para todos os tokens do prompt.
      \item \textbf{Por que mascarar o prompt?} Queremos que o gradiente venha apenas dos tokens que o modelo deve \emph{gerar}; gastar capacidade prevendo instruções (que o usuário escreve) não é o objetivo. É a prática padrão (Alpaca, LLaMA-Chat, HF TRL), embora haja evidência de que incluir a perda no prompt ajuda em alguns regimes (ex.: respostas curtas com prompts longos). Vale mencionar como decisão de projeto, não como lei.
      \item A figura ilustra isso: a caixa azul (instrução) contribui com gradiente zero; apenas a caixa verde (resposta) contribui para atualizar os pesos.
      \item \textbf{Hipótese do alinhamento superficial} (Zhou et al., LIMA): o conhecimento factual já está no modelo após o pré-treinamento. O SFT ensina apenas o \emph{formato} e o \emph{estilo} de resposta. Por isso, 1000 exemplos bem curados podem ser suficientes.
    \end{itemize}
  }
  \end{frame}

  %--- Slide: Datasets SFT ---
  \begin{frame}{Datasets para Instruction Tuning}
    \begin{columns}[T]
      \begin{column}{0.48\textwidth}
        \begin{block}{Datasets clássicos}
          \small
          \begin{itemize}
            \item \textbf{FLAN} \cite{chung2022flan}: 1.8k tarefas NLP;
              diversidade $>$ volume
            \item \textbf{ShareGPT}: conversas reais de usuários do ChatGPT;
              qualidade variável
            \item \textbf{LIMA} \cite{zhou2023lima}: 1.000 exemplos
              \textbf{muito bem curados} $\approx$ 52k do Alpaca
            \item \textbf{OpenHermes}: mistura de ShareGPT, GPT-4 e outros;
              qualidade alta e consistente
          \end{itemize}
        \end{block}
      \end{column}
      \begin{column}{0.48\textwidth}
        \begin{block}{Quantidade vs.\ Qualidade}
          \begin{center}
          \resizebox{0.95\linewidth}{!}{%
          \begin{tabular}{lcc}
            \toprule
            \textbf{Dataset} & \textbf{Exemplos} & \textbf{Qualidade} \\
            \midrule
            Alpaca (GPT-3.5) & 52\,000  & Baixa--Média \\
            ShareGPT         & 90\,000  & Média \\
            LIMA             & 1\,000   & \textcolor{DarkGreen}{Alta} \\
            OpenHermes-2.5   & 1\,000\,000 & Alta \\
            \bottomrule
          \end{tabular}}
          \end{center}
          \vspace{0.2em}
          \textcolor{DarkRed}{\small \textbf{Lição do LIMA:} ``1000 exemplos
          cuidadosamente selecionados superam 52000 gerados automaticamente.''}
        \end{block}
      \end{column}
    \end{columns}
  \note{
    \textbf{Pontos sobre datasets SFT:}
    \begin{itemize}
      \item \textbf{Alpaca} foi gerado por self-instruct via GPT-3.5: barato, mas com erros factuais e respostas genéricas. É um marco histórico (prova de conceito), não recomendado para produção.
      \item \textbf{LIMA insight:} a ``hipótese do alinhamento superficial'': o modelo já sabe \emph{o quê} dizer após o pré-treinamento; o SFT ensina apenas \emph{como} formatar a resposta. Por isso, 1000 exemplos excepcionalmente bons superam 52000 medíocres.
      \item \textbf{FLAN} demonstrou que a diversidade de tarefas (sumarização, Q\&A, tradução, etc.) melhora a generalização mesmo sem aumentar o volume de dados.
      \item Discussão útil: perguntar aos alunos como identificariam um exemplo de ``alta qualidade'' na prática (isso conecta com curadoria, avaliação humana e critérios de RLHF).
    \end{itemize}
  }
  \end{frame}

  %--- Slide: Chat templates ---
  \begin{frame}{\textit{Chat Templates} e Tokens Especiais}
    Modelos de chat usam \textbf{tokens especiais} para delimitar turnos da conversa, permitindo multi-turno e controle de papéis.
    Tokens comuns
          \begin{itemize}
            \item \texttt{<|im\_start|>}, \texttt{<|im\_end|>} (ChatML / Qwen)
            \item \texttt{[INST]}, \texttt{[/INST]} (LLaMA-2 chat)
            \item \texttt{<|start\_header\_id|>} (LLaMA-3)
            \item \texttt{<s>} (BOS), \texttt{</s>} (EOS)
          \end{itemize}
        
        O template define \textbf{onde calcular a perda}:
        apenas nos tokens gerados pelo assistente.
        \textcolor{DarkBlue}{\textbf{Exemplo: formato ChatML}}

        \centering
        \begin{tikzpicture}[font=\scriptsize\ttfamily]
          \node[draw=MidGray, fill=LightGray, rounded corners,
                text width=10.8cm, inner sep=5pt] {
            \textcolor{DarkBlue}{\textbf{<|im\_start|>system}}
            Você é um assistente útil.
            \textcolor{DarkBlue}{\textbf{<|im\_end|>}}\\[2pt]
            \textcolor{DarkGreen}{\textbf{<|im\_start|>user}}
            Explique redes neurais.
            \textcolor{DarkGreen}{\textbf{<|im\_end|>}}\\[2pt]
            \textcolor{DarkOrange}{\textbf{<|im\_start|>assistant}}
            \textbf{\textcolor{DarkRed}{[\textit{loss} calculada aqui]}}
            Redes neurais são\ldots
            \textcolor{DarkOrange}{\textbf{<|im\_end|>}}
          };
        \end{tikzpicture}
  \note{
    \textbf{Sobre a figura do formato ChatML:}
    \begin{itemize}
      \item Os tokens especiais (\texttt{<|im\_start|>}, \texttt{<|im\_end|>}) são adicionados ao vocabulário durante o treino do tokenizador. Eles têm IDs próprios e o modelo aprende a tratá-los como delimitadores de turno.
      \item A perda é calculada \textbf{apenas} nos tokens em laranja (resposta do assistente). Os tokens do sistema e do usuário são mascarados.
      \item \textbf{Consequência prática:} usar o template errado em inferência (ex.: usar o formato LLaMA-2 num modelo treinado com ChatML) degrada o desempenho significativamente: o modelo não ``reconhece'' onde começa sua resposta.
      \item Os diferentes templates refletem decisões de projeto de diferentes equipes (não há padronização universal, embora ChatML/OpenAI esteja se tornando referência).
    \end{itemize}
  }
  \end{frame}

  %--- Slide: Diferenças de otimização ---
  \begin{frame}{SFT vs.\ Pré-treinamento: Diferenças de Otimização}
    \begin{columns}[T]
      \begin{column}{0.48\textwidth}
        \begin{block}{Pré-treinamento}
          \begin{itemize}
            \item LR alto: $10^{-4}$ -- $3\times10^{-4}$
            \item Centenas de bilhões de tokens
            \item \textit{Batch size} muito grande (milhões de tokens)
            \item \textit{Weight decay}: 0.1
            \item Foco: \textbf{absorver conhecimento}
          \end{itemize}
        \end{block}
      \end{column}
      \begin{column}{0.48\textwidth}
        \begin{block}{Fine-tuning (SFT)}
          \begin{itemize}
            \item LR baixo: $10^{-6}$ -- $2\times10^{-5}$
            \item Milhares a milhões de exemplos
            \item \textit{Batch size} menor (128--512 sequências)
            \item Poucas épocas (1--3) para evitar \textit{overfitting}
            \item Foco: \textbf{adaptar comportamento}
          \end{itemize}
        \end{block}
      \end{column}
    \end{columns}

    \vspace{0.2em}
    \begin{alertblock}{Risco de \textit{Catastrophic Forgetting}}
      \small LR alto ou muitas épocas podem fazer o modelo ``esquecer'' conhecimento
      adquirido no pré-treinamento. Use LR baixo e monitore \textit{benchmarks}
      gerais ao longo do SFT.
    \end{alertblock}
  \end{frame}

  %======================================================================================
  % SECTION 3: RLHF COM PPO
  %======================================================================================
  \section{RLHF com PPO}

  %--- Slide: Pipeline RLHF ---
  \begin{frame}{Pipeline Completo do RLHF \cite{ouyang2022instructgpt,schulman2017ppo}}
    \centering
    \begin{tikzpicture}[
      node distance=0.4cm and 0.55cm,
      box/.style={draw, rounded corners, fill=CellBlue, text width=2.7cm,
                  align=center, font=\small, minimum height=1cm, inner sep=4pt},
      arr/.style={->, >=Stealth, thick}
    ]
      \node[box, fill=CellGreen] (sft)   {1.~\textbf{Modelo SFT}\\ (ponto de partida)};
      \node[box, fill=CellYellow, right=of sft] (rm)
            {2.~\textbf{Reward Model}\\ treinado em comparações};
      \node[box, fill=CellRed, right=of rm] (ppo)
            {3.~\textbf{PPO}\\ otimiza política usando RM};
      \node[box, fill=CellPurple, right=of ppo] (llm)
            {4.~\textbf{LLM alinhado}\\ (ex.: InstructGPT, Claude)};

      \draw[arr] (sft) -- (rm);
      \draw[arr] (rm)  -- (ppo);
      \draw[arr] (ppo) -- (llm);

      % KL penalty annotation
      \node[draw=DarkRed, fill=LightOrange, rounded corners, font=\scriptsize,
            text width=2.0cm, align=center, below=0.3cm of ppo] (kl)
            {Penalidade KL mantém política próxima do modelo SFT};
      \draw[->, DarkRed, dashed] (kl) -- (ppo);
      \draw[->, DarkRed, dashed] (kl.west) to[bend left=20] (sft.south);
    \end{tikzpicture}

    \vspace{0.2em}
    \begin{columns}
      \begin{column}{0.48\textwidth}
        \textbf{Etapa 1:} amostrar pares de respostas do modelo SFT\\
        \textbf{Etapa 2:} humanos escolhem a resposta preferida\\
        \textbf{Etapa 3:} treinar o RM com \textit{Bradley-Terry}
      \end{column}
      \begin{column}{0.48\textwidth}
        \textbf{Etapa 4:} usar PPO para otimizar a política,
        maximizando recompensa do RM com penalidade KL em relação ao SFT
      \end{column}
    \end{columns}

    \vspace{0.5em}
    \centering\textcolor{DarkRed}{\small O passo PPO mantém \textbf{4 modelos} em memória: (1) política $\pi_\theta$, (2) referência $\pi_{\text{ref}}$, (3) \textit{reward model} $r_\phi$, (4) \textit{value function} $V_\phi$, ou seja, ${\sim}4\times$ a VRAM do modelo base.}
  \note{
    \textbf{Sobre o diagrama do pipeline RLHF:}
    \begin{itemize}
      \item \textbf{Custo computacional:} durante o passo PPO, são necessários 4 modelos simultâneos em memória: (1) política sendo treinada $\pi_\theta$, (2) política de referência (SFT congelado) $\pi_{\text{ref}}$, (3) reward model $r_\phi$, (4) value function (crítico). Isso exige $\sim$4$\times$ a VRAM do modelo base.
      \item A seta tracejada vermelha (penalidade KL) conecta o PPO de volta ao SFT: enfatizar que o SFT é o ``ponto de ancoragem''. Sem ele, o PPO pode derivar para comportamentos degenerados.
      \item \textbf{InstructGPT} (Ouyang et al., 2022) foi o primeiro trabalho a demonstrar esse pipeline em escala. O ChatGPT usa variantes desse mesmo processo.
      \item O RM é treinado \emph{uma vez} (offline) com dados de preferência humana e depois \emph{congelado} durante o PPO. Atualizar o RM durante o PPO criaria um alvo móvel instável.
    \end{itemize}
  }
  \end{frame}

  %--- Slide: Dados de Preferência ---
  \begin{frame}{Dados de Preferência: Pares Escolhido/Rejeitado}
    \centering
    \textcolor{DarkBlue}{\textbf{Formato dos dados para RLHF, DPO e variantes}}
    \vspace{0.5em}

    \fcolorbox{MidGray}{LightGray}{\parbox{0.88\textwidth}{%
      \textcolor{DarkBlue}{\textbf{Prompt $x$:}}\enspace
      \textit{``Explique o que é aprendizado de máquina de forma simples.''}%
    }}

    \vspace{0.5em}
    \begin{columns}[T]
      \begin{column}{0.48\textwidth}
        \setbeamercolor{block title}{bg=DarkGreen!80!black, fg=white}
        \setbeamercolor{block body}{bg=CellGreen}
        \begin{block}{\textbf{Resposta escolhida $y_w$ (\checkmark)}}
          \textit{``Aprendizado de máquina é um método pelo qual computadores
          aprendem padrões a partir de dados, sem serem explicitamente
          programados para cada tarefa.''}
        \end{block}
        \vspace{0.3em}
        \centering\small Tripla: $(x,\; y_w,\; y_l)$, usada pelo RLHF (para treinar o RM), pelo DPO e variantes
      \end{column}
      \begin{column}{0.48\textwidth}
        \setbeamercolor{block title}{bg=DarkRed!80!black, fg=white}
        \setbeamercolor{block body}{bg=CellRed}
        \begin{block}{\textbf{Resposta rejeitada $y_l$ ($\times$)}}
          \textit{``ML é quando você faz o computador aprender coisas.
          É tipo uma coisa de dados. Basicamente é inteligência artificial
          mas diferente.''}
        \end{block}
        \vspace{0.3em}
        \centering\small\textcolor{MidGray}{HH-RLHF \cite{bai2022hhrlhf},
        UltraFeedback \cite{cui2024ultrafeedback}, OpenHermes}
      \end{column}
    \end{columns}
  \note{
    \textbf{Sobre os dados de preferência:}
    \begin{itemize}
      \item \textbf{Coleta:} um prompt $x$ é enviado ao modelo (ou a múltiplos modelos), gerando duas respostas. Anotadores humanos escolhem qual preferem (ou modelos GPT-4 atuam como juízes: RLAIF, Reinforcement Learning from AI Feedback).
      \item \textbf{Critérios de preferência:} utilidade, veracidade, ausência de conteúdo nocivo, clareza. Os critérios variam por dataset.
      \item \textbf{HH-RLHF} (Anthropic, 170k pares): foco em \textit{helpfulness} e \textit{harmlessness}. \textbf{UltraFeedback} (64k): gerado por GPT-4 como árbitro, multiaspecto. \textbf{OpenHermes}: instruções sintéticas de alta qualidade.
      \item A mesma tripla $(x, y_w, y_l)$ é usada por RLHF (para treinar o Reward Model) e por DPO, IPO, SimPO etc.\ (diretamente na loss de alinhamento). O GRPO \emph{não} usa pares: usa uma recompensa escalar por resposta (verificador ou RM).
    \end{itemize}
  }
  \end{frame}

  %--- Slide: Reward Model ---
  \begin{frame}{Treinamento do \textit{Reward Model}}
    \begin{columns}
      \begin{column}{0.52\textwidth}
        O \textit{Reward Model} (RM) é treinado para \textbf{prever preferência humana}
        entre pares de respostas para o mesmo prompt.

        \vspace{0.4em}
        \textbf{Coleta de dados:} para cada prompt $x$, amostram-se duas respostas
        $(y_w, y_l)$ onde $y_w$ é a preferida pelos anotadores.

        \vspace{0.4em}
        \begin{block}{Função de perda do RM (Bradley-Terry)}
          \small
          \[
            \mathcal{L}_{\text{RM}} = -\mathbb{E}_{(x,y_w,y_l)}\big[
              \log \sigma\!\left(r_\phi(x,y_w) - r_\phi(x,y_l)\right)
            \big]
          \]
          onde $r_\phi(x,y)$ é o \textit{scalar reward} predito.
        \end{block}

        \vspace{0.2em}
        {\small É a verossimilhança de um modelo de \textit{ranking}: equivale a \textbf{classificação binária} com \textit{logit} $= r_\phi(y_w){-}r_\phi(y_l)$.}
      \end{column}
      \begin{column}{0.45\textwidth}
        \textcolor{DarkBlue}{\textbf{Datasets de comparações humanas}}
        \vspace{0.3em}

        \resizebox{\linewidth}{!}{%
        \begin{tabular}{lcc}
          \toprule
          \textbf{Dataset} & \textbf{Pares} & \textbf{Fonte} \\
          \midrule
          HH-RLHF (Anthropic) & 170k  & Humanos \\
          WebGPT Comparisons  & 19k   & Humanos \\
          SHP                 & 385k  & Reddit \\
          UltraFeedback       & 64k   & GPT-4 \\
          \bottomrule
        \end{tabular}}

        \vspace{0.3em}
        \textcolor{MidGray}{\small O RM geralmente parte do \textbf{mesmo modelo SFT}
        com uma cabeça de regressão escalar adicionada.}
      \end{column}
    \end{columns}
  \note{
    \textbf{Intuição da perda Bradley-Terry:}
    \begin{itemize}
      \item O modelo de Bradley-Terry assume que a probabilidade de $y_w$ ser preferido sobre $y_l$ é $P(y_w \succ y_l) = \sigma(r(y_w) - r(y_l))$.
      \item A perda $-\log\sigma(\Delta r)$ é simplesmente a \textbf{log-verossimilhança negativa} desse modelo de ranking. Minimizá-la maximiza a diferença $r(y_w) - r(y_l)$, empurrando a recompensa da resposta preferida para cima e da rejeitada para baixo.
      \item \textbf{Analogia:} é idêntica à perda de classificação binária com logit $= r(y_w) - r(y_l)$. O RM aprende uma função de ``quanto melhor é $y_w$ que $y_l$?''.
      \item O RM inicializa \emph{do} modelo SFT (não do zero) porque já tem boa representação de linguagem. Apenas uma camada linear escalar é adicionada no topo.
      \item HH-RLHF da Anthropic contém pares helpfulness (escolha a resposta mais útil) e harmlessness (escolha a resposta mais segura): duas dimensões separadas.
    \end{itemize}
  }
  \end{frame}

  %--- Slide: Exemplo Numérico Reward Model ---
  \begin{frame}{Exemplo Numérico: \textit{Reward Model}}
    \textbf{Prompt $x$:} ``Explique o que é uma rede neural.'' \quad
    \textbf{Humanos preferem} $y_w$ (detalhada) sobre $y_l$ (vaga).

    \vspace{0.4em}
    \begin{columns}[T]
      \begin{column}{0.50\textwidth}
        \setbeamercolor{block title}{bg=DarkGreen!75!black, fg=white}
        \setbeamercolor{block body}{bg=CellGreen}
        \begin{block}{\textbf{Cenário A: RM bem treinado}}
          \small
          \begin{tabular}{lc}
            \toprule
            \textbf{Resposta} & $r_\phi(x,y)$ \\
            \midrule
            $y_w$ (detalhada) & $+2.1$ \\
            $y_l$ (vaga)      & $-0.8$ \\
            \bottomrule
          \end{tabular}

          \vspace{0.3em}
          \begin{align*}
            \Delta r &= (+2.1) - (-0.8) = \mathbf{+2.9} \\[1pt]
            \sigma(\Delta r) &\approx 0.948 \\[1pt]
            \mathcal{L}_{\text{RM}} &= -\log(0.948) \approx \mathbf{0.053}
          \end{align*}
          \centering \textcolor{DarkGreen}{\checkmark\ RM concorda com o rótulo humano}
        \end{block}
      \end{column}
      \begin{column}{0.50\textwidth}
        \setbeamercolor{block title}{bg=DarkRed!75!black, fg=white}
        \setbeamercolor{block body}{bg=CellRed}
        \begin{block}{\textbf{Cenário B: RM mal treinado}}
          \small
          \begin{tabular}{lc}
            \toprule
            \textbf{Resposta} & $r_\phi(x,y)$ \\
            \midrule
            $y_w$ (detalhada) & $-0.8$ \\
            $y_l$ (vaga)      & $+2.1$ \\
            \bottomrule
          \end{tabular}

          \vspace{0.3em}
          \begin{align*}
            \Delta r &= (-0.8) - (+2.1) = \mathbf{-2.9} \\[1pt]
            \sigma(\Delta r) &\approx 0.052 \\[1pt]
            \mathcal{L}_{\text{RM}} &= -\log(0.052) \approx \mathbf{2.96}
          \end{align*}
          \centering \textcolor{DarkRed}{$\times$\ RM discorda; gradiente grande}
        \end{block}
      \end{column}
    \end{columns}
  \note{
    \textbf{Leitura do exemplo:}
    \begin{itemize}
      \item O rótulo humano é fixo nos dois cenários: $y_w$ é sempre a resposta preferida. O que muda é o estado do RM, não os dados.
      \item \textbf{Cenário A (bem treinado):} o RM aprendeu a ordenar corretamente, atribuindo reward maior à resposta detalhada. $\Delta r = +2.9$ leva a $\sigma(\Delta r) \approx 0.95$ e loss pequena ($\approx 0.05$), sinalizando que pouco ajuste é necessário.
      \item \textbf{Cenário B (mal treinado):} o RM ainda confunde as respostas, dando reward maior à vaga. $\Delta r = -2.9$ leva a $\sigma(\Delta r) \approx 0.05$ e loss $\approx 2.96$, quase $60\times$ maior. Esse gradiente grande empurra $r_\phi(x,y_w)$ para cima e $r_\phi(x,y_l)$ para baixo, corrigindo a ordenação.
      \item A curva sigmoide da perda Bradley-Terry satura: quando $|\Delta r|$ é grande e na direção certa, o gradiente é próximo de zero (modelo já convencido); na direção errada, o gradiente é máximo em torno de $\Delta r = 0$.
    \end{itemize}
  }
  \end{frame}

  %--- Slide: Value function e vantagem ---
  \begin{frame}{O Quarto Modelo: \textit{Value Function} e Vantagem}
    \small
    \begin{columns}[T]
      \begin{column}{0.46\textwidth}
        \begin{block}{O crítico $V_\phi$ opera por token}
          Para cada prefixo $s_t = (x, y_{<t})$, estima a recompensa futura: $V_\phi(s_t) \approx \mathbb{E}\big[R \mid s_t\big]$.
          Mesma arquitetura da política com uma cabeça escalar, treinado por regressão.
        \end{block}
        \begin{block}{Vantagem por token}
          \vspace{-0.6em}
          \[
            A_t = r_t + V_\phi(s_{t+1}) - V_\phi(s_t)
          \]
          \vspace{-1.1em}

          ``Quanto o token $t$ mudou a expectativa de recompensa?''
          No RLHF, $r_t = 0$ exceto no último token, que recebe a nota do RM.
        \end{block}
      \end{column}
      \begin{column}{0.52\textwidth}
        \textcolor{DarkBlue}{\textbf{Exemplo:}} $x$ = ``Quanto é $7\times 8$?'', RM dá $R = 0{,}2$
        \vspace{0.2em}

        \resizebox{\linewidth}{!}{%
        \begin{tabular}{clcccc}
          \toprule
          $t$ & token $y_t$ & $V_\phi(s_t)$ & $r_t$ & $V_\phi(s_{t+1})$ & $A_t$ \\
          \midrule
          1 & \texttt{É}     & $1{,}2$ & $0$     & $1{,}3$ & $\textcolor{DarkGreen}{+0{,}1}$ \\
          2 & \texttt{igual} & $1{,}3$ & $0$     & $1{,}4$ & $\textcolor{DarkGreen}{+0{,}1}$ \\
          3 & \texttt{a}     & $1{,}4$ & $0$     & $1{,}5$ & $\textcolor{DarkGreen}{+0{,}1}$ \\
          4 & \texttt{54}    & $1{,}5$ & $0{,}2$ & $0$ (fim) & $\textcolor{DarkRed}{-1{,}3}$ \\
          \bottomrule
        \end{tabular}}

        \vspace{0.3em}
        Os três primeiros tokens elevaram a expectativa e são levemente reforçados; o número errado recebe \textbf{toda a culpa}.

        \vspace{0.1em}
        \textcolor{MidGray}{Com uma \textit{baseline} só do prompt ($b = 1{,}2$), os quatro tokens receberiam a mesma vantagem $R - b = -1{,}0$.}

        \vspace{0.1em}
        \textcolor{DarkRed}{\textbf{Custo:} uma rede do tamanho da política, treinada junto com ela: o quarto modelo em memória.}
      \end{column}
    \end{columns}
  \note{
    \textbf{Por que o PPO precisa de um crítico:}
    \begin{itemize}
      \item \textbf{Gradiente de política (REINFORCE):} $\nabla J = \mathbb{E}[\nabla\log\pi_\theta(y_t|s_t)\,R]$. Funciona, mas a variância é enorme: toda ação de uma trajetória com $R>0$ é reforçada, mesmo as ruins. Subtrair qualquer função $b(s_t)$ que não dependa da ação mantém o gradiente esperado ($\mathbb{E}[\nabla\log\pi \cdot b] = b\,\nabla\sum_a \pi = 0$) e reduz a variância. A melhor baseline é aproximadamente $V(s_t)$.
      \item \textbf{Por token:} no RLHF a recompensa é esparsa (o RM só avalia a resposta completa, no último token), e a penalidade KL entra por token. O crítico é o que distribui esse sinal pelos tokens intermediários (\textit{credit assignment}): no exemplo, ele ``sabia'' que o prefixo ``É igual a'' estava indo bem ($V$ subindo) e que o ``54'' estragou tudo ($V$ cai de 1,5 para 0 e só entram 0,2 de reward).
      \item \textbf{Sobre a fórmula do slide:} $A_t = r_t + V(s_{t+1}) - V(s_t)$ é a vantagem de um passo (erro de TD, $\gamma = 1$). O PPO usa GAE (\textit{Generalized Advantage Estimation}, Schulman et al., 2016), que combina esse estimador com o retorno completo $R - V(s_t)$ via um parâmetro $\lambda$; no RLHF costuma-se usar $\gamma = 1$ e $\lambda = 0{,}95$, ou seja, quase o retorno completo, mas com a correção por token do crítico. Com $\lambda = 1$ o exemplo daria $A_t = R - V(s_t)$: $-1{,}0$, $-1{,}1$, $-1{,}2$, $-1{,}3$ (todos negativos, sem distinguir os tokens bons).
      \item \textbf{Comparação com a baseline do prompt:} a linha cinza é exatamente o que o GRPO faz (vantagem igual para todos os tokens da resposta, calculada a partir de uma baseline que não depende do prefixo). Funciona bem quando as respostas do grupo diferem por inteiro; perde a granularidade por token.
      \item \textbf{Inicialização:} o crítico normalmente parte dos pesos do Reward Model (já aprendeu a avaliar respostas) e é treinado com a perda quadrática $(V_\phi(s_t) - R_t)^2$ em paralelo ao clip surrogate da política. Ignoramos a penalidade KL nos números do exemplo para não poluir a tabela.
    \end{itemize}
  }
  \end{frame}

  %--- Slide: PPO objetivo ---
  \begin{frame}{Função Objetivo do PPO com Penalidade KL \cite{schulman2017ppo}}
    \begin{columns}
      \begin{column}{0.52\textwidth}
        O PPO otimiza a política $\pi_\theta$ para maximizar recompensa
        enquanto permanece próximo da política SFT de referência $\pi_{\text{ref}}$.

        \begin{block}{Objetivo RLHF-PPO}
          \[
            J(\theta) = \mathbb{E}_{x,y\sim\pi_\theta}\!\left[
              r_\phi(x,y) - \beta\, \text{KL}\!\left(\pi_\theta \Vert \pi_{\text{ref}}\right)
            \right]
          \]
        \end{block}
        \vspace{-0.2em}
        {\scriptsize\textcolor{MidGray}{Na prática a KL vira uma \textbf{penalidade por token}: subtrai-se $\beta\log\frac{\pi_\theta(t)}{\pi_{\text{ref}}(t)}$ do \textit{reward} de cada token.}}

        \begin{block}{\textit{Clip Surrogate} do PPO}
          \[
            L^{\text{CLIP}} = \mathbb{E}_t\!\left[
              \min\!\left(r_t A_t,\;
              \text{clip}(r_t, 1{-}\varepsilon, 1{+}\varepsilon) A_t
              \right)
            \right]
          \]
          onde $r_t = \pi_\theta / \pi_{\text{old}}$ e $A_t$ é a vantagem.
        \end{block}
      \end{column}
      \begin{column}{0.45\textwidth}
        \textcolor{DarkBlue}{\textbf{Papel do coeficiente KL} $\beta$}
        \vspace{0.1em}

        \resizebox{\linewidth}{!}{%
        \begin{tabular}{lp{3.5cm}}
          \toprule
          $\beta$ pequeno & Liberdade para maximizar reward; risco de \textit{reward hacking} \\
          \midrule
          $\beta$ grande  & Permanece próximo do SFT; pode não melhorar o alinhamento \\
          \bottomrule
        \end{tabular}}

        \textcolor{DarkRed}{\textbf{Instabilidades típicas:}}
        \begin{itemize}
          \item \textit{Reward hacking}: modelo ``engana'' o RM
          \item \textit{KL explosion}: divergência excessiva da política
        \end{itemize}
      \end{column}
    \end{columns}
  \note{
    \textbf{Intuição das duas equações do PPO:}
    \begin{itemize}
      \item \textbf{Objetivo RLHF-PPO:} o termo $-\beta\,\text{KL}(\pi_\theta \Vert \pi_{\text{ref}})$ penaliza o modelo por se afastar do SFT. Na prática, isso é implementado como uma \emph{penalidade por token}: para cada token gerado, subtrai-se $\beta \log(\pi_\theta(t)/\pi_{\text{ref}}(t))$ do reward (o reward do RM entra só no último token da resposta).
      \item \textbf{Clip Surrogate:} $r_t = \pi_\theta(a_t|s_t)/\pi_{\text{old}}(a_t|s_t)$ é a razão de importância. O clipping em $[1-\varepsilon, 1+\varepsilon]$ (tipicamente $\varepsilon=0.2$) impede atualizações muito grandes: se $r_t > 1+\varepsilon$, o ganho adicional não é considerado. Isso mantém a política próxima de $\pi_{\text{old}}$ a cada iteração.
      \item $A_t$ é a \textbf{vantagem}: quanto melhor foi a ação $a_t$ em relação à baseline (valor esperado). Se $A_t > 0$, queremos aumentar $r_t$; se $A_t < 0$, queremos diminuir.
      \item \textbf{Reward hacking:} exemplo clássico: modelos PPO aprendem a gerar respostas longas com repetição, pois o RM mal treinado dá rewards maiores a respostas longas. A penalidade KL deveria prevenir isso, mas $\beta$ precisa ser calibrado.
    \end{itemize}
  }
  \end{frame}

  %======================================================================================
  % SECTION 4: DPO
  %======================================================================================
  \section{Direct Preference Optimization (DPO)}

  %--- Slide: Motivação DPO ---
  \begin{frame}{DPO: Motivação}
    \begin{columns}
      \begin{column}{0.52\textwidth}
        \textbf{Problema com RLHF-PPO:}
        \begin{itemize}
          \item Requer treinamento de um \textit{Reward Model} separado
          \item Otimização RL é instável e sensível a hiperparâmetros
          \item Precisa manter 4 modelos em memória simultaneamente
            (política, referência, RM, \textit{value function})
        \end{itemize}

        \textbf{Ideia do DPO} (Rafailov et al., 2023): \textit{eliminar o RM, reparametrizando o problema diretamente na política.}

        \textcolor{DarkGreen}{\small O DPO usa a forma da solução ótima do RLHF para dispensar o RM.}
      \end{column}
      \begin{column}{0.45\textwidth}
        \textcolor{DarkBlue}{\textbf{Comparação de infraestrutura}}
        \vspace{0.3em}

        \resizebox{\linewidth}{!}{%
        \begin{tabular}{lcc}
          \toprule
          \textbf{Componente} & \textbf{RLHF} & \textbf{DPO} \\
          \midrule
          Reward Model        & \textcolor{DarkRed}{\checkmark} & \textcolor{DarkGreen}{\texttimes} \\
          Value Function      & \textcolor{DarkRed}{\checkmark} & \textcolor{DarkGreen}{\texttimes} \\
          Política de referência & \checkmark & \checkmark \\
          Geração online      & \textcolor{DarkRed}{\checkmark} & \textcolor{DarkGreen}{\texttimes} \\
          Dados de comparação & \checkmark & \checkmark \\
          \bottomrule
        \end{tabular}}
      \end{column}
    \end{columns}

    \vspace{0.2em}
    \centering\textcolor{MidGray}{\small \textit{Online} $=$ \textbf{on-policy} (amostras frescas); DPO é \textbf{off-policy} (pares pré-coletados).}
  \note{
    \textbf{Sobre a tabela de infraestrutura DPO vs.\ RLHF:}
    \begin{itemize}
      \item A tabela resume o custo de engenharia: RLHF-PPO precisa de RM treinado (pipeline separado), value function (outro modelo do mesmo tamanho), geração online (lenta, custosa). DPO elimina tudo isso.
      \item \textbf{Geração online vs.\ offline:} PPO gera amostras frescas a cada iteração (on-policy), garantindo que o RM avalie respostas da política atual. DPO usa dados de comparação coletados previamente (off-policy), o que pode ser uma limitação se a política mudar muito.
      \item A ``solução ótima fechada'' mencionada no slide é a chave do DPO: da teoria de otimização com restrição KL, sabe-se que $\pi^*(y|x) \propto \pi_{\text{ref}}(y|x)\exp(r(x,y)/\beta)$. DPO inverte essa relação para expressar $r$ em função de $\pi$.
    \end{itemize}
  }
  \end{frame}

  %--- Slide: DPO enunciado ---
  \begin{frame}{A Perda do DPO \cite{rafailov2023dpo}}
    \small
    A política ótima do RLHF com penalidade KL tem a forma $\pi^*(y|x) \propto \pi_{\text{ref}}(y|x)\exp\!\big(r(x,y)/\beta\big)$.
    Invertendo-a, o \textit{reward} vira uma função da política; no Bradley-Terry a constante de normalização cancela e sobra:
    \begin{block}{Perda DPO}
      \vspace{-0.4em}
      \[
        \mathcal{L}_{\text{DPO}}(\theta) = -\mathbb{E}_{(x,y_w,y_l)}\!\left[\log\sigma\!\left(
          \beta\log\frac{\pi_\theta(y_w|x)}{\pi_{\text{ref}}(y_w|x)}
          - \beta\log\frac{\pi_\theta(y_l|x)}{\pi_{\text{ref}}(y_l|x)}
        \right)\right]
      \]
      \vspace{-0.6em}
    \end{block}
    \begin{columns}[T]
      \begin{column}{0.50\textwidth}
        \textcolor{DarkBlue}{\textbf{Como ler}}
        \begin{itemize}
          \item É a perda do Reward Model com o \textit{reward} trocado pelo \textbf{\textit{reward} implícito} $\hat r_\theta(x,y) = \beta\log\frac{\pi_\theta(y|x)}{\pi_{\text{ref}}(y|x)}$
          \item O modelo de linguagem \emph{é} o reward model
          \item $\beta$ controla o afastamento de $\pi_{\text{ref}}$ (típico: 0,1 a 0,5)
        \end{itemize}
      \end{column}
      \begin{column}{0.47\textwidth}
        \textcolor{DarkGreen}{\textbf{O que o gradiente faz}}
        \begin{itemize}
          \item Aumenta $\log\pi_\theta(y_w|x)$ e diminui $\log\pi_\theta(y_l|x)$
          \item Peso $\sigma\big(\hat r_\theta(y_l) - \hat r_\theta(y_w)\big)$: pares ainda mal ordenados dominam a atualização
          \item Só precisa de \textit{log-probs} nos pares: sem geração, sem RM, sem crítico
        \end{itemize}
      \end{column}
    \end{columns}
  \note{
    \textbf{Sobre a perda DPO (sem a derivação completa):}
    \begin{itemize}
      \item \textbf{Esboço da derivação, se perguntarem:} (1) a solução do RLHF com KL é $\pi^* = \pi_{\text{ref}}\exp(r/\beta)/Z(x)$, com $Z(x)$ a função de partição; (2) isolando, $r = \beta\log(\pi^*/\pi_{\text{ref}}) + \beta\log Z(x)$; (3) no Bradley-Terry só entra a diferença $r(y_w) - r(y_l)$, e $\beta\log Z(x)$ cancela; (4) trocando $\pi^*$ por $\pi_\theta$ e maximizando a verossimilhança dos pares, sai a perda. A parte intratável ($Z$) era irrelevante para preferências. O artigo original (Rafailov et al., 2023, apêndice A) tem a derivação completa em uma página.
      \item \textbf{Gradiente:} $\nabla_\theta\mathcal{L} = -\beta\,\mathbb{E}[\sigma(\hat r_\theta(y_l) - \hat r_\theta(y_w))\,(\nabla\log\pi_\theta(y_w|x) - \nabla\log\pi_\theta(y_l|x))]$. O peso é a probabilidade que o RM implícito atribui à ordem errada: pares já bem ordenados quase não contribuem.
      \item \textbf{Implementação:} $\log\pi_\theta(y|x) = \sum_t\log\pi_\theta(y_t|x,y_{<t})$ sai de um forward com \textit{teacher forcing}; os log-probs de $\pi_{\text{ref}}$ são constantes e podem ser pré-computados. Dois forwards por par. Por isso o DPO cabe numa GPU onde o PPO não cabe.
      \item \textbf{Off-policy:} os pares vêm de $\pi_{\text{ref}}$ ou de outros modelos, não de $\pi_\theta$. Se $\pi_\theta$ se afastar muito, os pares deixam de ser representativos. É a principal fraqueza frente ao PPO e a motivação das variantes \textit{online}.
      \item \textbf{Fenômeno comum:} $\log\pi_\theta(y_w|x)$ frequentemente \emph{cai} durante o treino, apenas menos que $\log\pi_\theta(y_l|x)$; a perda só olha para a diferença. Variantes como IPO e SimPO tentam mitigar isso.
    \end{itemize}
  }
  \end{frame}

  %======================================================================================
  % SECTION 5: OUTROS MÉTODOS DE ALINHAMENTO
  %======================================================================================
  \section{Outros Métodos de Alinhamento}

  %--- Slide: GRPO ---
  \begin{frame}{GRPO: \textit{Group Relative Policy Optimization} \cite{shao2024deepseekmath,deepseekr1}}
        Proposto no DeepSeekMath (Shao et al., 2024) e popularizado pelo DeepSeek-R1, o GRPO elimina a
        \textit{value function} do PPO usando como \textbf{baseline a média de um grupo}: para cada prompt, amostram-se $G$ respostas $y_1,\ldots,y_G$ com recompensas $r_1,\ldots,r_G$.
        \vspace{-0.2em}
          \[
            J_{\text{GRPO}} = \mathbb{E}\!\bigg[\frac{1}{G}\sum_{i=1}^{G}
              \min\!\left(\rho_i\hat{A}_i,\;\text{clip}(\rho_i,1{-}\varepsilon,1{+}\varepsilon)\,\hat{A}_i\right)
              \;-\; \beta\,\text{KL}\!\left(\pi_\theta \Vert \pi_{\text{ref}}\right)
            \bigg]
          \]
          onde $\rho_i = \dfrac{\pi_\theta(y_i|x)}{\pi_{\text{old}}(y_i|x)}$ é a razão de importância e
          $\hat{A}_i = \dfrac{r_i - \text{mean}(\{r_j\})}{\text{std}(\{r_j\})}$ é a vantagem relativa ao grupo.

        \begin{itemize}\small
          \item Sem \textit{value function}: um modelo a menos em memória (3 em vez de 4)
          \item Vantagem por $z$-score no grupo: sem variância no grupo, $\hat{A}_i{=}0$ e não há gradiente
          \item Funciona bem com \textit{reward} esparso e verificável (resposta final certa ou errada)
          \item O termo KL continua ancorando a política ao modelo de referência
        \end{itemize}

        \textcolor{DarkGreen}{\footnotesize \textbf{Resultado:} DeepSeek-R1 alcança
        desempenho competitivo com OpenAI o1 em AIME 2024 e MATH-500 usando GRPO
        com verificação automática de respostas.}
  \note{
    \textbf{Intuição da equação GRPO:}
    \begin{itemize}
      \item \textbf{Atribuição:} o GRPO aparece no DeepSeekMath (fev.\ 2024), onde foi usado para matemática. O R1 (jan.\ 2025) mostrou que o mesmo algoritmo, com reward verificável, faz emergir raciocínio longo. Citar os dois.
      \item \textbf{Notação:} $\rho_i$ é a mesma razão de importância do PPO (lá chamada de $r_t$); usamos $\rho$ aqui porque $r_i$ já denota a recompensa. No paper, a razão e o clip são calculados \emph{por token} e a média é feita sobre os tokens de cada resposta; o slide mostra a versão por sequência para simplificar.
      \item \textbf{Estimativa de vantagem por grupo:} a vantagem de cada resposta é sua recompensa normalizada pelo grupo: $\hat{A}_i = (r_i - \bar{r})/\sigma_r$. Isso é simplesmente um \emph{z-score} (quanto melhor que a média do grupo?).
      \item Por que isso substitui a value function? A value function do PPO estima $E[r | s_t]$ para cada estado intermediário. O GRPO substitui essa estimativa por uma baseline simples: a média das recompensas do grupo para o mesmo prompt. Funciona bem quando o reward é \emph{esparso} (só recebido no final).
      \item \textbf{KL:} o GRPO original mantém a penalidade KL (com um estimador não enviesado $\frac{\pi_{\text{ref}}}{\pi_\theta} - \log\frac{\pi_{\text{ref}}}{\pi_\theta} - 1$), somada diretamente à perda em vez de embutida no reward. Trabalhos posteriores (DAPO, Dr.\ GRPO) removem ou modificam esse termo.
      \item O reward no DeepSeek-R1 é verificado automaticamente: para matemática, compara-se a resposta final com a solução correta via parser simbólico. Para código, rodam-se os testes. Isso elimina a necessidade de um RM neural.
    \end{itemize}
  }
  \end{frame}

  %--- Slide: Exemplo Numérico GRPO ---
  \begin{frame}{Exemplo Numérico: GRPO}
    \textbf{Problema:} ``Quanto é $12 \times 7$?'' \quad (resposta correta: $84$) \quad
    Grupo de $G = 4$ respostas amostradas.

    \vspace{0.3em}
    \begin{columns}[T]
      \begin{column}{0.50\textwidth}
        \resizebox{\linewidth}{!}{%
        \begin{tabular}{clcc}
          \toprule
          $i$ & Resposta $y_i$ & $r_i$ & $\hat{A}_i$ \\
          \midrule
          1 & ``84'' (correto) & 1 & \textcolor{DarkGreen}{$+1.0$} \\
          2 & ``82'' (errado)  & 0 & \textcolor{DarkRed}{$-1.0$} \\
          3 & ``84'' (correto) & 1 & \textcolor{DarkGreen}{$+1.0$} \\
          4 & ``76'' (errado)  & 0 & \textcolor{DarkRed}{$-1.0$} \\
          \bottomrule
        \end{tabular}}

        \vspace{0.4em}
        \begin{align*}
          \overline{r} &= \tfrac{1+0+1+0}{4} = 0.5, \qquad \sigma_r = 0.5 \\[4pt]
          \hat{A}_i &= \frac{r_i - 0.5}{0.5}
        \end{align*}
      \end{column}
      \begin{column}{0.47\textwidth}
        \begin{block}{Efeito na política}
          \begin{itemize}
            \item $\hat{A}_i > 0$: resposta \textbf{melhor que a média}
              $\Rightarrow$ política \textcolor{DarkGreen}{reforçada}
            \item $\hat{A}_i < 0$: resposta \textbf{pior que a média}
              $\Rightarrow$ política \textcolor{DarkRed}{penalizada}
          \end{itemize}
        \end{block}

        \vspace{0.4em}
        \textcolor{MidGray}{\small Sem \textit{value function}: o baseline é
        a \textbf{média das recompensas do grupo}.
        Com recompensas binárias e metade do grupo acertando,
        $\hat{A}_i \in \{-1,\,+1\}$.}
      \end{column}
    \end{columns}
  \note{
    \textbf{Intuição do exemplo:}
    \begin{itemize}
      \item $y_1$ e $y_3$ acertam ($r=1$) e recebem vantagem $\hat{A} = +1$: o gradiente empurra a política para produzir mais essas respostas.
      \item $y_2$ e $y_4$ erram ($r=0$) e recebem $\hat{A} = -1$: o gradiente afasta a política delas.
      \item Se \emph{todas} acertassem, $\hat{A}_i = 0$ para todas: sem gradiente, pois o modelo já é ótimo para esse grupo. Isso mostra que o GRPO só aprende quando há variância no grupo.
      \item O z-score normaliza a vantagem independentemente da escala do reward. Para recompensas contínuas (ex.\ notas 0--10), o comportamento é o mesmo.
    \end{itemize}
  }
  \end{frame}

  %--- Slide: RLVR e modelos de raciocínio ---
  \begin{frame}{RL com Recompensas Verificáveis (RLVR) \cite{deepseekr1,lambert2024tulu3}}
    \small
    \begin{columns}[T]
      \begin{column}{0.50\textwidth}
        \begin{block}{Ideia}
          Trocar o \textit{reward model} neural por um \textbf{verificador}: gabarito em matemática, testes unitários em código, checagem de formato. Recompensa barata, exata e difícil de enganar.
        \end{block}
        \begin{block}{Receita do DeepSeek-R1}
          \begin{enumerate}
            \item \textit{Cold start}: SFT curto com cadeias de raciocínio
            \item GRPO com \textit{reward} de acurácia + formato
            \item \textit{Rejection sampling} das melhores gerações e novo SFT
            \item Nova rodada de RL, agora com preferências humanas também
          \end{enumerate}
        \end{block}
      \end{column}
      \begin{column}{0.47\textwidth}
        \textcolor{DarkBlue}{\textbf{O que emerge durante o RL}}
        \begin{itemize}
          \item Respostas ficam \textbf{mais longas}: o modelo aprende a ``pensar'' antes de responder
          \item Autoverificação e retrocesso (``espera, vou conferir\ldots'') sem supervisão direta
        \end{itemize}
        \vspace{0.2em}
        \textcolor{DarkRed}{\textbf{Riscos e limites}}
        \begin{itemize}
          \item \textit{Reward hacking} no verificador (formato do gabarito, testes fracos)
          \item Só funciona onde há verificador: escrita e diálogo seguem com RM ou pares
          \item Custo de inferência cresce com cadeias longas
        \end{itemize}
      \end{column}
    \end{columns}
  \note{
    \textbf{RLVR e raciocínio:}
    \begin{itemize}
      \item O termo RLVR (\textit{Reinforcement Learning with Verifiable Rewards}) foi cunhado no Tülu 3 (Lambert et al., 2024); a ideia é a mesma do reward baseado em resultado do DeepSeekMath e do R1.
      \item \textbf{R1-Zero vs.\ R1:} o R1-Zero fez RL direto sobre o modelo base, sem SFT, e obteve raciocínio forte mas com texto ilegível e mistura de idiomas. O R1 adiciona o cold start e as rodadas de rejection sampling para recuperar legibilidade e utilidade geral.
      \item \textbf{Comprimento crescente:} a curva de tokens por resposta ao longo do treino RL é a figura mais citada do paper. Ligar com o custo de inferência: raciocínio longo é \textit{test-time compute}.
      \item \textbf{Reward hacking:} verificadores fracos (comparação de string, testes incompletos) são explorados rapidamente. Boas práticas: verificadores simbólicos (SymPy), múltiplos testes ocultos, penalidade de formato.
      \item Enfatizar a limitação: sem verificador, voltamos ao RM neural ou a pares de preferência. Modelos de fronteira combinam os dois regimes no mesmo pós-treino.
    \end{itemize}
  }
  \end{frame}

  %--- Slide: SimPO ---
  \begin{frame}{SimPO: \textit{Reward} Sem Modelo de Referência \cite{meng2024simpo}}
    \textbf{Duas críticas ao DPO:} o \textit{reward} implícito $\beta\log\frac{\pi_\theta(y|x)}{\pi_{\text{ref}}(y|x)}$ não é o critério usado na \textbf{geração}, que olha só para $\pi_\theta$; e manter $\pi_{\text{ref}}$ em memória custa uma cópia inteira do modelo.
    \begin{block}{Perda SimPO}
      \vspace{-0.4em}
      \[
        \mathcal{L}_{\text{SimPO}} = -\log\sigma\!\bigg(
          \frac{\beta}{|y_w|}\log\pi_\theta(y_w|x)
          - \frac{\beta}{|y_l|}\log\pi_\theta(y_l|x) - \gamma \bigg)
      \]
    \end{block}
    \begin{columns}[T]
      \begin{column}{0.52\textwidth}
        \begin{itemize}\small
          \item \textit{Reward} $=$ log-verossimilhança \textbf{média por token}: alinhado com a geração e sem viés de comprimento
          \item Margem $\gamma > 0$ exige separação mínima entre $y_w$ e $y_l$
          \item Sem $\pi_{\text{ref}}$: $\sim$20\% menos memória e treino mais rápido
          \item \textcolor{DarkRed}{Sem âncora, $\gamma$ e $\beta$ exigem ajuste fino}
        \end{itemize}
      \end{column}
      \begin{column}{0.45\textwidth}
        \textcolor{DarkBlue}{\textbf{\textit{Reward} implícito de cada método}}
        \vspace{0.2em}

        \resizebox{\linewidth}{!}{%
        \begin{tabular}{lcc}
          \toprule
          \textbf{Método} & $r_\theta(x,y)$ & $\pi_{\text{ref}}$? \\
          \midrule
          DPO   & $\beta\log\dfrac{\pi_\theta(y|x)}{\pi_{\text{ref}}(y|x)}$ & sim \\[1.2em]
          SimPO & $\dfrac{\beta}{|y|}\log\pi_\theta(y|x)$ & não \\
          \bottomrule
        \end{tabular}}
      \end{column}
    \end{columns}
  \note{
    \textbf{Intuição da perda SimPO:}
    \begin{itemize}
      \item \textbf{Desalinhamento treino/geração no DPO:} o DPO otimiza $\log\pi_\theta - \log\pi_{\text{ref}}$, mas na hora de gerar usamos apenas $\pi_\theta$. Um par pode ter reward implícito correto ($y_w$ acima de $y_l$) e ainda assim $\pi_\theta(y_w) < \pi_\theta(y_l)$, se a referência já preferia muito $y_l$. SimPO remove essa discrepância.
      \item \textbf{Normalização por comprimento:} $\log\pi_\theta(y|x) = \sum_{t}\log\pi_\theta(y_t|x,y_{<t})$ é uma soma sobre os tokens. Sequências mais longas têm log-probabilidade mais negativa simplesmente por ter mais termos: dividir por $|y|$ normaliza isso, evitando que o modelo prefira respostas curtas (ou longas) apenas por viés de comprimento.
      \item \textbf{Margem $\gamma$:} garante que $\text{reward}(y_w) - \text{reward}(y_l) > \gamma$ (não apenas $> 0$). Isso previne que o modelo estagne com separação mínima entre chosen e rejected. Valores típicos: $\beta \in [2, 10]$, $\gamma/\beta \in [0.3, 1.5]$.
      \item \textbf{Sem âncora:} o preço de remover $\pi_{\text{ref}}$ é perder a garantia de que a política não se afasta demais do SFT. Na prática, SimPO usa poucas épocas e LR baixo para compensar.
    \end{itemize}
  }
  \end{frame}

  %--- Slide: Tabela comparativa ---
  \begin{frame}{Tabela Comparativa dos Métodos de Alinhamento}
    \centering
    \resizebox{0.95\linewidth}{!}{%
    \begin{tabular}{lccccc}
      \toprule
      \textbf{Método} &
      \textbf{Precisa de RM} &
      \textbf{Precisa de pares} &
      \textbf{\textit{Online}} &
      \textbf{Sem ref.} &
      \textbf{Caso de uso típico} \\
      \midrule
      \rowcolor{CellYellow}
      RLHF-PPO  & \textcolor{DarkGreen}{\checkmark} & \textcolor{DarkGreen}{\checkmark} &
        \textcolor{DarkGreen}{\checkmark} & \texttimes &
        Alinhamento geral (InstructGPT) \\
      DPO  & \texttimes & \textcolor{DarkGreen}{\checkmark} & \texttimes & \texttimes &
        Alinhamento simples, eficiente \\
      \rowcolor{CellBlue}
      SimPO & \texttimes & \textcolor{DarkGreen}{\checkmark} & \texttimes &
        \textcolor{DarkGreen}{\checkmark} &
        Sem dependência de referência \\
      GRPO & \textcolor{DarkGreen}{\checkmark}* & \texttimes & \textcolor{DarkGreen}{\checkmark} &
        \texttimes & Raciocínio matemático/código \\
      \bottomrule
      \multicolumn{6}{l}{\small * GRPO usa uma recompensa escalar por resposta: um verificador automático (RLVR) ou, quando não há verificador, um RM neural.}
    \end{tabular}}

    \vspace{0.5em}
    \begin{block}{Panorama em 2026}
      \textbf{SFT + DPO} (ou variantes \textit{offline}) segue sendo o caminho mais barato e o padrão em modelos abertos menores.
      Modelos de fronteira combinam, em várias rodadas, \textit{rejection sampling}, DPO \textit{online} e
      \textbf{RL com recompensas verificáveis} (GRPO e sucessores) para raciocínio, código e agentes.
      PPO clássico com \textit{value function} ficou restrito a equipes com infraestrutura dedicada.
    \end{block}
  \note{
    \textbf{Como usar esta tabela em aula:}
    \begin{itemize}
      \item Percorrer as colunas: ``Precisa de RM?'' elimina todos exceto RLHF-PPO e GRPO*. ``Online?'' elimina todos exceto RLHF-PPO e GRPO.
      \item A nota de rodapé do GRPO é importante: na maioria dos usos, o ``RM'' do GRPO é um verificador simbólico/algorítmico (solução matemática, execução de código), não uma rede neural treinada em preferências humanas. É muito mais barato e confiável. Nada impede, porém, usar GRPO com um RM neural em tarefas sem verificador.
      \item \textbf{Decisão prática para os alunos:} ``Você tem pares de preferência humana? Use DPO. Você quer alinhamento sem modelo de referência? Use SimPO. Você quer treinar raciocínio matemático/código com verificação automática? Use GRPO. Você precisa de alinhamento de última geração com orçamento computacional alto? Use RL online (PPO ou GRPO) com RM e verificadores combinados.''
      \item \textbf{Exemplos de pipelines reais:} Llama 3 (Grattafiori et al., 2024) fez seis rodadas de rejection sampling + SFT + DPO; Tülu 3 fez SFT + DPO + RLVR; DeepSeek-R1 e Qwen3 fazem SFT curto + GRPO em larga escala. A mensagem: não há um método vencedor, há uma composição.
    \end{itemize}
  }
  \end{frame}

  %--- Slide: Resumo do pipeline ---
  \begin{frame}{Resumo: Pipeline Completo de Treinamento de LLMs}
    \centering
    \begin{tikzpicture}[
      node distance=0.3cm and 0.45cm,
      stage/.style={draw, rounded corners=4pt, text width=1.8cm, align=center,
                    font=\small, minimum height=1.2cm, inner sep=4pt},
      arr/.style={->, >=Stealth, thick, DarkBlue}
    ]
      \node[stage, fill=CellGreen]   (pt)   {\textbf{Pré-treino}\\{\scriptsize Next-token pred.}\\{\scriptsize 1T--15T tokens}};
      \node[stage, fill=CellYellow,  right=of pt]  (sft)
            {\textbf{SFT}\\{\scriptsize Instruction tuning}\\{\scriptsize 1k--1M exemplos}};
      \node[stage, fill=CellBlue,    right=of sft] (rm)
            {\textbf{RM}\\{\scriptsize Pares preferidos}\\{\scriptsize (opcional)}};
      \node[stage, fill=CellRed,     right=of rm]  (align)
            {\textbf{Alinhamento}\\{\scriptsize PPO / DPO / SimPO}\\{\scriptsize GRPO (RLVR)}};
      \node[stage, fill=CellPurple,  right=of align] (deploy)
            {\textbf{Deploy}\\{\scriptsize Quantização, GGUF}\\{\scriptsize vLLM, TGI}};

      \draw[arr] (pt) -- (sft);
      \draw[arr] (sft) -- (rm);
      \draw[arr] (rm)  -- (align);
      \draw[arr] (align) -- (deploy);

      % Optional RM bypass for DPO
      \draw[arr, DarkGreen, dashed, bend right=30]
            (sft.south) to node[below, font=\scriptsize, DarkGreen]
            {DPO/SimPO (sem RM)} (align.south);
    \end{tikzpicture}

    \vspace{0.6em}
    \begin{columns}
      \begin{column}{0.32\textwidth}
        \begin{block}{\small Pré-treino}
          \scriptsize
          GPT-3, LLaMA, Falcon\\
          Datasets: The Pile, C4\\
          Precisão Mista, Paralelismo
        \end{block}
      \end{column}
      \begin{column}{0.32\textwidth}
        \begin{block}{\small SFT + Alinhamento}
          \scriptsize
          InstructGPT, ChatGPT\\
          LLaMA-Chat, Mistral-Instruct\\
          Qwen-Chat, DeepSeek-R1
        \end{block}
      \end{column}
      \begin{column}{0.32\textwidth}
        \begin{block}{\small Pesquisa ativa}
          \scriptsize
          RL com recompensas verificáveis\\
          Variantes do GRPO (DAPO, Dr.\ GRPO)\\
          \textit{Constitutional AI} / RLAIF
        \end{block}
      \end{column}
    \end{columns}
  \note{
    \textbf{Sobre o diagrama do pipeline completo:}
    \begin{itemize}
      \item A seta verde tracejada (``DPO/SimPO sem RM'') é a inovação central dos últimos 2 anos: pula o passo de treinar o RM, indo diretamente de SFT para alinhamento.
      \item \textbf{Deploy:} ferramentas como vLLM e TGI permitem servir modelos quantizados em produção. Mencionar que o deploy eficiente é tão importante quanto o treino: um modelo de 70B sem quantização não roda em hardware acessível.
      \item \textbf{Constitutional AI} (Anthropic): o modelo usa um conjunto de princípios (a ``Constituição'') para autoavaliar e revisar suas próprias respostas, gerando dados de preferência sem anotadores humanos na escala principal.
      \item Usar este diagrama como revisão ao final da aula: perguntar aos alunos onde cada método estudado se encaixa no pipeline.
    \end{itemize}
  }
  \end{frame}

  %======================================================================================
  % SECTION 6: OTIMIZAÇÕES DE TREINAMENTO
  %======================================================================================
  \section{Otimizações de Treinamento}

  %--- Slide: Motivação Precisão Mista ---
  \begin{frame}{Precisão Mista: Por que Usar Menos Bits? \cite{micikevicius2018mixed}}
    \begin{columns}[T]
      \begin{column}{0.50\textwidth}
        Treinar em \textbf{precisão total} (fp32) consome 16 bytes
        por parâmetro (pesos + gradientes + estados Adam).

        \begin{block}{Motivação}
          Reduzir a precisão de fp32 para fp16/bf16 durante as
          operações principais permite:
          \begin{itemize}
            \item \textbf{2$\times$ menos memória} para pesos e ativações
            \item \textbf{2--16$\times$ throughput} em Tensor Cores modernos
            \item Treinamento de modelos maiores com o mesmo hardware
          \end{itemize}
        \end{block}

        \textcolor{DarkRed}{\footnotesize \textbf{Atenção:} precisão insuficiente causa \textit{overflow/underflow} e instabilidade.}
      \end{column}
      \begin{column}{0.47\textwidth}
        \textcolor{DarkBlue}{\textbf{Throughput de GEMM (denso)}}
        \vspace{0.3em}

        \resizebox{\linewidth}{!}{%
        \begin{tabular}{lcc}
          \toprule
          \textbf{GPU} & \textbf{FP32 (CUDA cores)} & \textbf{BF16/FP16 (Tensor Cores)} \\
          \midrule
          A100 SXM & 19.5 TFLOPS & 312 TFLOPS \\
          H100 SXM & 67 TFLOPS   & 989 TFLOPS \\
          RTX 4090 & 83 TFLOPS   & 165 TFLOPS \\
          \bottomrule
          \multicolumn{3}{l}{\scriptsize GEMM = multiplicação de matrizes. Valores sem \textit{sparsity} estruturada.}
        \end{tabular}}

        \vspace{0.4em}
        \textcolor{MidGray}{\small \textbf{Tensor Cores} (GPUs Volta+) são
        unidades de hardware dedicadas à multiplicação de matrizes em
        fp16/bf16 (e TF32/FP8 nas gerações recentes), daí o salto de throughput acima.}
      \end{column}
    \end{columns}
  \note{
    \textbf{Por que GPUs são muito mais rápidas em fp16/bf16?}
    \begin{itemize}
      \item GPUs modernas (Volta+) possuem \emph{Tensor Cores}: unidades de hardware dedicadas a multiplicações de matrizes em fp16/bf16. Um A100 faz 19.5 TFLOPS em fp32 nos CUDA cores, mas 312 TFLOPS em bf16 nos Tensor Cores (16$\times$). No H100: 67 vs.\ 989 TFLOPS, e 1\,979 TFLOPS em FP8. As fichas técnicas da NVIDIA costumam anunciar o dobro desses valores ``com sparsity'' (2:4 estruturada), que não se aplica a treino comum; por isso a tabela usa os valores densos.
      \item Em FP32 também existe o modo \emph{TF32} dos Tensor Cores (19 bits: 8 de expoente, 10 de mantissa), ativado por padrão no PyTorch para convoluções e opcionalmente para matmul: 156 TFLOPS no A100, 495 no H100. É o motivo de ``fp32'' em GPU moderna já ser mais rápido do que parece.
      \item O bottleneck de memória também melhora: pesos em fp16 ocupam metade da memória, logo mais pesos cabem no cache L2/SRAM da GPU, reduzindo acessos à HBM.
      \item \textbf{Por que não treinar sempre em fp16?} O intervalo dinâmico do fp16 (máximo $\approx 65504$) causa \emph{overflow} de gradientes em modelos grandes. O bf16 tem o mesmo intervalo do fp32 (8 bits de expoente) e é mais robusto numericamente.
    \end{itemize}
  }
  \end{frame}

  %--- Slide: Representações de Ponto Flutuante ---
  \begin{frame}{Representações de Ponto Flutuante em Deep Learning}
    \begin{columns}[T]
      \begin{column}{0.60\textwidth}
        \centering
        \textcolor{DarkBlue}{\small\textbf{Layout de bits \quad
          \textcolor{DarkRed}{$\blacksquare$ Sinal} \quad
          \textcolor{DarkBlue}{$\blacksquare$ Expoente} \quad
          \textcolor{DarkGreen}{$\blacksquare$ Mantissa}}}
        \vspace{0.3em}

        \begin{tikzpicture}
          \pgfmathsetmacro{\bw}{0.19}
          \pgfmathsetmacro{\bh}{0.58}
          \pgfmathsetmacro{\sx}{0.10}

          % Macro: draw one row of bit-fields
          % #1=y-center, #2=label, #3=S, #4=E, #5=M bits, #6=total label
          % Approach: use \fill+\draw for rectangles (pure path, no node props)

          % ---- FP32 (1+8+23 = 32 bits) ----
          \pgfmathsetmacro{\y}{3.80}
          \node[anchor=east, font=\footnotesize\bfseries] at (-0.05, \y) {FP32};
          \fill[PastelRed]  (\sx,            {\y-\bh/2}) rectangle ++({\bw*1},  \bh);
          \draw[DarkRed!60, thin] (\sx, {\y-\bh/2}) rectangle ++({\bw*1}, \bh);
          \fill[CellBlue]   ({\sx+\bw*1},    {\y-\bh/2}) rectangle ++({\bw*8},  \bh);
          \draw[DarkBlue!60, thin] ({\sx+\bw*1}, {\y-\bh/2}) rectangle ++({\bw*8}, \bh);
          \fill[PastelGreen]({\sx+\bw*9},    {\y-\bh/2}) rectangle ++({\bw*23}, \bh);
          \draw[DarkGreen!60, thin] ({\sx+\bw*9}, {\y-\bh/2}) rectangle ++({\bw*23}, \bh);
          \node[font=\tiny\bfseries, text=DarkBlue]  at ({\sx+\bw*5},    \y) {8};
          \node[font=\tiny\bfseries, text=DarkGreen] at ({\sx+\bw*20.5}, \y) {23};
          \node[anchor=west, font=\scriptsize, text=MidGray] at ({\sx+\bw*32+0.08}, \y) {32 bits};

          % ---- FP16 (1+5+10 = 16 bits) ----
          \pgfmathsetmacro{\y}{2.95}
          \node[anchor=east, font=\footnotesize\bfseries] at (-0.05, \y) {FP16};
          \fill[PastelRed]  (\sx,         {\y-\bh/2}) rectangle ++({\bw*1},  \bh);
          \draw[DarkRed!60, thin] (\sx, {\y-\bh/2}) rectangle ++({\bw*1}, \bh);
          \fill[CellBlue]   ({\sx+\bw*1}, {\y-\bh/2}) rectangle ++({\bw*5},  \bh);
          \draw[DarkBlue!60, thin] ({\sx+\bw*1}, {\y-\bh/2}) rectangle ++({\bw*5}, \bh);
          \fill[PastelGreen]({\sx+\bw*6}, {\y-\bh/2}) rectangle ++({\bw*10}, \bh);
          \draw[DarkGreen!60, thin] ({\sx+\bw*6}, {\y-\bh/2}) rectangle ++({\bw*10}, \bh);
          \node[font=\tiny\bfseries, text=DarkBlue]  at ({\sx+\bw*3.5}, \y) {5};
          \node[font=\tiny\bfseries, text=DarkGreen] at ({\sx+\bw*11},  \y) {10};
          \node[anchor=west, font=\scriptsize, text=MidGray] at ({\sx+\bw*32+0.08}, \y) {16 bits};

          % ---- BF16 (1+8+7 = 16 bits) ----
          \pgfmathsetmacro{\y}{2.10}
          \node[anchor=east, font=\footnotesize\bfseries] at (-0.05, \y) {BF16};
          \fill[PastelRed]  (\sx,         {\y-\bh/2}) rectangle ++({\bw*1}, \bh);
          \draw[DarkRed!60, thin] (\sx, {\y-\bh/2}) rectangle ++({\bw*1}, \bh);
          \fill[CellBlue]   ({\sx+\bw*1}, {\y-\bh/2}) rectangle ++({\bw*8}, \bh);
          \draw[DarkBlue!60, thin] ({\sx+\bw*1}, {\y-\bh/2}) rectangle ++({\bw*8}, \bh);
          \fill[PastelGreen]({\sx+\bw*9}, {\y-\bh/2}) rectangle ++({\bw*7}, \bh);
          \draw[DarkGreen!60, thin] ({\sx+\bw*9}, {\y-\bh/2}) rectangle ++({\bw*7}, \bh);
          \node[font=\tiny\bfseries, text=DarkBlue]  at ({\sx+\bw*5},    \y) {8};
          \node[font=\tiny\bfseries, text=DarkGreen] at ({\sx+\bw*12.5}, \y) {7};
          \node[anchor=west, font=\scriptsize, text=MidGray] at ({\sx+\bw*32+0.08}, \y) {16 bits};

          % ---- FP8 E4M3 (1+4+3 = 8 bits) ----
          \pgfmathsetmacro{\y}{1.25}
          \node[anchor=east, font=\tiny\bfseries, align=right] at (-0.05, \y) {FP8\\E4M3};
          \fill[PastelRed]  (\sx,         {\y-\bh/2}) rectangle ++({\bw*1}, \bh);
          \draw[DarkRed!60, thin] (\sx, {\y-\bh/2}) rectangle ++({\bw*1}, \bh);
          \fill[CellBlue]   ({\sx+\bw*1}, {\y-\bh/2}) rectangle ++({\bw*4}, \bh);
          \draw[DarkBlue!60, thin] ({\sx+\bw*1}, {\y-\bh/2}) rectangle ++({\bw*4}, \bh);
          \fill[PastelGreen]({\sx+\bw*5}, {\y-\bh/2}) rectangle ++({\bw*3}, \bh);
          \draw[DarkGreen!60, thin] ({\sx+\bw*5}, {\y-\bh/2}) rectangle ++({\bw*3}, \bh);
          \node[font=\tiny\bfseries, text=DarkBlue]  at ({\sx+\bw*3},   \y) {4};
          \node[font=\tiny\bfseries, text=DarkGreen] at ({\sx+\bw*6.5}, \y) {3};
          \node[anchor=west, font=\scriptsize, text=MidGray] at ({\sx+\bw*32+0.08}, \y) {8 bits};

          % ---- FP8 E5M2 (1+5+2 = 8 bits) ----
          \pgfmathsetmacro{\y}{0.40}
          \node[anchor=east, font=\tiny\bfseries, align=right] at (-0.05, \y) {FP8\\E5M2};
          \fill[PastelRed]  (\sx,         {\y-\bh/2}) rectangle ++({\bw*1}, \bh);
          \draw[DarkRed!60, thin] (\sx, {\y-\bh/2}) rectangle ++({\bw*1}, \bh);
          \fill[CellBlue]   ({\sx+\bw*1}, {\y-\bh/2}) rectangle ++({\bw*5}, \bh);
          \draw[DarkBlue!60, thin] ({\sx+\bw*1}, {\y-\bh/2}) rectangle ++({\bw*5}, \bh);
          \fill[PastelGreen]({\sx+\bw*6}, {\y-\bh/2}) rectangle ++({\bw*2}, \bh);
          \draw[DarkGreen!60, thin] ({\sx+\bw*6}, {\y-\bh/2}) rectangle ++({\bw*2}, \bh);
          \node[font=\tiny\bfseries, text=DarkBlue]  at ({\sx+\bw*3.5}, \y) {5};
          \node[font=\tiny\bfseries, text=DarkGreen] at ({\sx+\bw*7},   \y) {2};
          \node[anchor=west, font=\scriptsize, text=MidGray] at ({\sx+\bw*32+0.08}, \y) {8 bits};
        \end{tikzpicture}
      \end{column}
      \begin{column}{0.37\textwidth}
        \textcolor{DarkBlue}{\textbf{Impacto dos campos}}
        \vspace{0.2em}

        {\footnotesize
        \begin{itemize}\setlength{\itemsep}{1pt}
          \item \textbf{Expoente}: \textbf{intervalo} dinâmico
          \item \textbf{Mantissa}: \textbf{precisão} entre valores
        \end{itemize}}

        \vspace{0.3em}
        \textcolor{DarkBlue}{\textbf{Alcance representável (valores positivos)}}

        \vspace{0.2em}
        \resizebox{\linewidth}{!}{%
        \begin{tabular}{lcc}
          \toprule
          \textbf{Formato} & \textbf{Mín.\ normal} & \textbf{Máx.\ finito} \\
          \midrule
          FP32      & $\approx 1.18 \times 10^{-38}$ & $\approx 3.4 \times 10^{38}$ \\
          FP16      & $\approx 6.1 \times 10^{-5}$   & $65\,504$ \\
          BF16      & $\approx 1.18 \times 10^{-38}$ & $\approx 3.4 \times 10^{38}$ \\
          FP8 E4M3  & $\approx 1.56 \times 10^{-2}$  & $448$ \\
          FP8 E5M2  & $\approx 6.1 \times 10^{-5}$   & $57\,344$ \\
          \bottomrule
        \end{tabular}}
      \end{column}
    \end{columns}
        \begin{block}{\small Comparação chave}
        \small
          \textbf{BF16} = mesmo expoente do FP32 (8 bits), evita overflow de
          gradientes; \textbf{FP16} = mais mantissa (10 bits) mas intervalo
          de $\sim$9 ordens de grandeza, contra $\sim$76 do FP32.
          \textbf{FP8 E5M2} (mais range): gradientes. \textbf{FP8 E4M3} (mais precisão): pesos/ativações.
        \end{block}
  \note{
    \textbf{Como ler o diagrama de bits:}
    \begin{itemize}
      \item Valor representado: $(-1)^S \times 2^{E - \text{bias}} \times (1 + M/2^{|M|})$. O expoente tem ``bias'' subtraído (ex.: FP32 bias = 127; o menor campo de expoente normal é $E=1$, que representa $2^{1-127} = 2^{-126} \approx 1.18\times10^{-38}$; $E=0$ é reservado para zero e subnormais, $E=255$ para $\pm\infty$ e NaN).
      \item \textbf{FP32 vs BF16}: ambos têm 8 bits de expoente (mesmo intervalo dinâmico $[\sim\!3.4\times10^{-38}, \sim\!3.4\times10^{38}]$). BF16 sacrifica 16 bits de mantissa (23→7) para caber em 16 bits. Para gradientes de redes neurais, a precisão de 7 bits de mantissa geralmente é suficiente.
      \item \textbf{FP16}: 5 bits de expoente $\to$ máximo $\approx 65504$. Gradientes de modelos grandes frequentemente excedem isso, causando \textit{inf/NaN}. Por isso o AMP usa \textit{loss scaling} com FP16 mas não com BF16.
      \item \textbf{FP8}: introduzido nas GPUs H100 (Hopper). E4M3 para forward (mais precisão); E5M2 para backward (mais range). Throughput $\approx 2\times$ BF16 em H100.
    \end{itemize}
  }
  \end{frame}

  %--- Slide: AMP ---
  \begin{frame}{Precisão Mista Automática (AMP)}
    \small
    \textbf{Pesos mestres:} cópia dos parâmetros em fp32 usada pelo otimizador.
    \textit{Forward} e \textit{backward} usam uma cópia em fp16/bf16 (rápida nos \textit{Tensor Cores}),
    mas o AdamW acumula na cópia mestre em fp32 para não descartar atualizações minúsculas.

    \begin{columns}[T]
      \begin{column}{0.48\textwidth}
        \centering
        \begin{tikzpicture}[
          node distance=0.3cm,
          box/.style={draw, rounded corners=3pt, minimum height=0.45cm,
                     text width=3.4cm, align=center, font=\scriptsize,
                     inner sep=2pt},
          arr/.style={->, >=Stealth, thick},
        ]
          \node[box, fill=CellBlue] (w32)
            {\textbf{Pesos mestres} (fp32)};
          \node[box, fill=PastelGreen!70, below=of w32] (cast)
            {\textbf{Cast} fp32 $\to$ fp16/bf16};
          \node[box, fill=PastelGreen!70, below=of cast] (fwd)
            {\textbf{Forward pass} (fp16/bf16)};
          \node[box, fill=PastelYellow, below=of fwd] (loss)
            {\textbf{Loss} + \textit{loss scaling} $\times S$};
          \node[box, fill=PastelGreen!70, below=of loss] (bwd)
            {\textbf{Backward} (fp16/bf16 grads)};
          \node[box, fill=PastelRed!70, below=of bwd] (unscale)
            {\textbf{Unscale} ($\div S$) + clip grads};
          \node[box, fill=CellBlue, below=of unscale] (upd)
            {\textbf{Atualizar pesos} (AdamW fp32)};

          \draw[arr] (w32) -- (cast);
          \draw[arr] (cast) -- (fwd);
          \draw[arr] (fwd) -- (loss);
          \draw[arr] (loss) -- (bwd);
          \draw[arr] (bwd) -- (unscale);
          \draw[arr] (unscale) -- (upd);
          \draw[arr] (upd.east) -- ++(0.4,0) |- (w32.east);
        \end{tikzpicture}
      \end{column}
      \begin{column}{0.49\textwidth}
        \begin{block}{\textit{Loss Scaling} (necessário com FP16)}
          Gradientes em fp16 sofrem \textbf{underflow} (valores pequenos
          $\to$ zero). Solução: escalar a loss:
          \[
            \tilde{\mathcal{L}} = S \cdot \mathcal{L},\quad
            g_{\text{fp32}} = (\tilde{g}_{\text{fp16}}) / S
          \]
          $S$ é ajustado dinamicamente: sobe se sem overflow,
          desce se overflow detectado.
        \end{block}

        \begin{alertblock}{BF16 dispensa \textit{loss scaling}}
          \footnotesize Mesmo intervalo dinâmico do FP32 $\Rightarrow$ sem risco
          de overflow em gradientes. \textbf{Preferir BF16} em GPUs Ampere+.
        \end{alertblock}
      \end{column}
    \end{columns}
  \note{
    \textbf{Sobre AMP e loss scaling:}
    \begin{itemize}
      \item A cópia mestre em fp32 existe porque as atualizações do AdamW acumulam pequenas diferenças: a resolução da mantissa do fp16 (10 bits $\approx 3$ dígitos decimais) não é suficiente para acumular gradientes minúsculos sem perda de informação.
      \item \textbf{GradScaler dinâmico (PyTorch):} começa com $S = 2^{16}$. Se ocorrer overflow/NaN no gradiente: $S \leftarrow S / 2$ e o passo é descartado. Após 2\,000 passos consecutivos sem overflow (\texttt{growth\_interval}): $S \leftarrow S \times 2$. Assim $S$ fica sempre perto do maior valor que ainda não estoura.
      \item O backward em fp16/bf16 é possível porque os GEMMs do backward pass (multiplicações de ativações e erros) se beneficiam dos Tensor Cores: só as somas de gradientes e as atualizações do otimizador precisam de fp32.
    \end{itemize}
  }
  \end{frame}

  %--- Slide: Paralelismo e memória de ativações ---
  \begin{frame}{Treinamento Distribuído: Paralelismo e Memória \cite{rajbhandari2020zero,shoeybi2019megatron,chen2016checkpoint}}
    \small
    Um modelo de 70B em bf16 ocupa 140 GB só em pesos: \textbf{nenhuma GPU} o contém. O treino distribui o trabalho em três eixos.
    \begin{columns}[T]
      \begin{column}{0.62\textwidth}
        \resizebox{\linewidth}{!}{%
        \begin{tabular}{lp{4.6cm}p{3.6cm}}
          \toprule
          \textbf{Eixo} & \textbf{O que se divide} & \textbf{Custo de comunicação} \\
          \midrule
          \textit{Data parallel} (DP) & O \textit{batch}: cada GPU tem uma réplica completa & \textit{All-reduce} dos gradientes a cada passo \\
          \rowcolor{LightGray}
          \textit{Tensor parallel} (TP) & Cada matriz de pesos, dentro da camada (Megatron-LM) & \textit{All-reduce} em toda camada, exige NVLink \\
          \textit{Pipeline parallel} (PP) & As camadas, em estágios sequenciais (GPipe) & Só ativações na fronteira; sofre com ``bolhas'' \\
          \rowcolor{LightGray}
          ZeRO / FSDP & Os \textbf{estados do otimizador, gradientes e pesos} entre as réplicas de DP & \textit{All-gather} dos pesos antes de cada camada \\
          \bottomrule
        \end{tabular}}

        \vspace{0.3em}
        \textcolor{DarkBlue}{\textbf{Na prática (LLaMA, Megatron):} TP dentro do nó (8 GPUs), PP entre nós, DP com ZeRO por cima; $\text{TP}\times\text{PP}\times\text{DP}$ é o total de GPUs.}
      \end{column}
      \begin{column}{0.36\textwidth}
        \begin{block}{\textit{Gradient checkpointing}}
          As ativações do \textit{forward} dominam a memória em sequências longas.
          Guardar só as de algumas camadas e \textbf{recomputar} o resto no \textit{backward}
          troca $\sim$30\% de computação por memória $\mathcal{O}(\sqrt{L})$ em vez de $\mathcal{O}(L)$.
        \end{block}
        \textcolor{MidGray}{ZeRO-3/FSDP e \textit{checkpointing} permitem \textit{full fine-tuning} de 7B--70B em poucos nós; para uma GPU só, a resposta é PEFT.}
      \end{column}
    \end{columns}
  \note{
    \textbf{Paralelismo e memória:}
    \begin{itemize}
      \item \textbf{DP puro} replica tudo: com AdamW em precisão mista, cada GPU precisa de 16 bytes/parâmetro, o que limita DP a modelos de $\sim$1B em GPUs de 24 GB. Por isso DP quase sempre vem com ZeRO.
      \item \textbf{ZeRO} (Rajbhandari et al., 2020) tem três estágios: ZeRO-1 divide os estados do otimizador (12 dos 16 bytes), ZeRO-2 também os gradientes, ZeRO-3 também os pesos. Com ZeRO-3 a memória por GPU cai para $16/N_{\text{DP}}$ bytes/parâmetro mais ativações. FSDP do PyTorch é a implementação nativa do ZeRO-3.
      \item \textbf{TP} (Megatron-LM) divide cada matriz $\mathbf{W}$ por colunas ou linhas entre GPUs, de modo que a atenção e a MLP de uma camada sejam calculadas em paralelo. Exige comunicação a cada camada, por isso fica confinado às GPUs de um mesmo nó (NVLink, $\sim$600--900 GB/s).
      \item \textbf{PP} coloca camadas diferentes em GPUs diferentes; o problema são as bolhas (GPUs ociosas esperando o estágio anterior). Mitiga-se dividindo o batch em micro-batches (GPipe, 1F1B).
      \item \textbf{Gradient checkpointing} (Chen et al., 2016): sem ele, um transformador guarda várias ativações intermediárias por camada e por token (entrada da atenção, scores, saída da MLP etc.). Com checkpoints a cada $\sqrt{L}$ camadas, o custo é um forward extra ($\sim$1/3 do passo). É o padrão em HF \texttt{gradient\_checkpointing=True}.
      \item Ligar com a próxima seção: tudo isso é infraestrutura de quem treina o modelo base ou faz full fine-tuning. Para a maioria dos alunos, o caminho realista é PEFT em uma GPU.
    \end{itemize}
  }
  \end{frame}

  %======================================================================================
  % SECTION 7: AJUSTE EFICIENTE EM PARÂMETROS (PEFT)
  %======================================================================================
  \section{Ajuste Eficiente em Parâmetros (PEFT)}

  %--- Slide: Motivação PEFT ---
  \begin{frame}{Motivação: Custo do \textit{Fine-tuning} Completo}
    \begin{columns}[T]
      \begin{column}{0.50\textwidth}
        Fazer \textit{fine-tuning} completo de um LLM exige armazenar em GPU:
        \begin{itemize}
          \item \textbf{Pesos}: 2 bytes/parâmetro (bf16)
          \item \textbf{Gradientes}: 2 bytes/parâmetro
          \item \textbf{Otimizador} (AdamW): 12 bytes/parâmetro (cópia mestre fp32 + $m$ + $v$)
          \item \textbf{Ativações}: dependem de \textit{batch} e comprimento
        \end{itemize}

        \vspace{0.2em}
        \begin{alertblock}{Exemplo: modelo de 7B parâmetros}
          Pesos + gradientes + otimizador $\approx$ \textbf{112 GB}:
          fora do alcance de GPUs de consumo (RTX~4090: 24~GB).
        \end{alertblock}
      \end{column}
      \begin{column}{0.47\textwidth}
        \textcolor{DarkBlue}{\textbf{Memória por parâmetro}}
        \vspace{0.3em}

        \resizebox{\linewidth}{!}{%
        \begin{tabular}{lcc}
          \toprule
          \textbf{Componente} & \textbf{fp32} & \textbf{fp16/bf16} \\
          \midrule
          Pesos               & 4 bytes & 2 bytes \\
          Gradientes          & 4 bytes & 2 bytes \\
          Cópia mestre (fp32) & ---     & 4 bytes \\
          Adam $m$ (1\textordfeminine\ mom.)  & 4 bytes & 4 bytes \\
          Adam $v$ (2\textordfeminine\ mom.)  & 4 bytes & 4 bytes \\
          \midrule
          \textbf{Total}      & \textbf{16 bytes} & \textbf{16 bytes} \\
          \bottomrule
          \multicolumn{3}{l}{\scriptsize Em precisão mista, a cópia mestre e os momentos do Adam ficam em fp32.}
        \end{tabular}}

        \vspace{0.4em}
        \textcolor{DarkGreen}{\small \textbf{Solução PEFT:} treinar apenas uma
        fração pequena dos parâmetros, mantendo o modelo base congelado.}
      \end{column}
    \end{columns}
  \note{
    \textbf{Sobre a tabela de memória e o custo do fine-tuning completo:}
    \begin{itemize}
      \item \textbf{Por que os estados Adam ficam em fp32 mesmo com pesos em fp16?} Mixed-precision training: os pesos são mantidos em fp16 para a forward pass (economia de memória), mas uma cópia ``mestre'' em fp32 é mantida para as atualizações (precisão numérica). Os momentos do Adam ($m$, $v$) também ficam em fp32. Resultado: $2 + 2 + 4 + 4 + 4 = 16$ bytes/param.
      \item Para 7B params: $7\times10^9 \times 16 = 112$ GB (confirma o número no slide). Uma A100 80GB não chega para fine-tuning completo de um modelo 7B!
      \item \textbf{Ativações} não estão na tabela porque dependem de batch e comprimento; \textit{gradient checkpointing} (visto no slide de treinamento distribuído) reduz esse custo de $O(L)$ para $O(\sqrt{L})$.
      \item A PEFT resolve o problema dos gradientes e estados do otimizador: se só 0,1\% dos params são treináveis, os gradientes e o estado Adam de apenas 0,1\% dos params precisam ser armazenados.
    \end{itemize}
  }
  \end{frame}

  %--- Slide: LoRA ---
  \begin{frame}{LoRA: \textit{Low-Rank Adaptation} \cite{hu2022lora}}
    \begin{columns}[T]
      \begin{column}{0.52\textwidth}
        Hu et al.\ (2022) decompõem a atualização de pesos como um produto
        \textbf{\textit{low rank}}:

        \begin{block}{Decomposição LoRA}
          Para uma camada $\mathbf{W}_0 \in \mathbb{R}^{d \times k}$ congelada,
          adicionamos um desvio treinável:
          \[
            \mathbf{h} = \mathbf{W}_0 \mathbf{x} + \tfrac{\alpha}{r}\, \mathbf{B} \mathbf{A} \mathbf{x}
          \]
          onde $\mathbf{B} \in \mathbb{R}^{d \times r}$, $\mathbf{A} \in \mathbb{R}^{r \times k}$,
          $r \ll \min(d,k)$ e $\alpha$ é um fator de escala.
        \end{block}

        \begin{itemize}
          \item $\mathbf{A}$ gaussiana, $\mathbf{B}$ com \textbf{zeros}: $\Delta\mathbf{W}{=}\mathbf{0}$ no início, o modelo parte do base
          \item Somente $\mathbf{A}$ e $\mathbf{B}$ recebem gradiente
          \item Parâmetros treináveis: $r(d{+}k) \ll dk$
        \end{itemize}
      \end{column}
      \begin{column}{0.45\textwidth}
        \textcolor{DarkBlue}{\textbf{Configurações típicas}}
        \vspace{0.3em}

        \resizebox{\linewidth}{!}{%
        \begin{tabular}{lcc}
          \toprule
          \textbf{Hiperparâmetro} & \textbf{Valor típico} & \textbf{Notas} \\
          \midrule
          \textit{rank} $r$     & 4--64      & Menor = menos parâmetros \\
          $\alpha$               & $r$--$2r$  & Controla a escala \\
          \textit{lora\_dropout} & 0.05       & Regularização \\
          Camadas alvo           & Q, V       & Mínimo; ver abaixo \\
          \bottomrule
        \end{tabular}}

        \vspace{0.4em}
        \textcolor{DarkRed}{\small \textbf{Exemplo:} LLaMA-7B com $r{=}16$
        em Q e V: $\sim$8,4M parâmetros treináveis ($\sim$0,12\% do total);
        a memória de gradientes e otimizador praticamente desaparece.}

        \vspace{0.3em}
        \textcolor{MidGray}{\small Adaptar todas as projeções lineares
        (Q, K, V, O, \textit{gate}, \textit{up}, \textit{down}) melhora o desempenho
        ao custo de mais parâmetros; é o padrão atual (QLoRA).}
      \end{column}
    \end{columns}
  \note{
    \textbf{Intuição da fórmula LoRA ($\mathbf{h} = \mathbf{W}_0 \mathbf{x} + \frac{\alpha}{r}\mathbf{B}\mathbf{A}\mathbf{x}$):}
    \begin{itemize}
      \item $\Delta\mathbf{W} = \frac{\alpha}{r}\mathbf{B}\mathbf{A}$ é a \emph{atualização} aprendida. O rank de $\Delta\mathbf{W}$ é no máximo $r \ll d, k$, daí o nome \textit{Low-Rank Adaptation}. A hipótese do paper: a atualização necessária para adaptar um modelo pré-treinado tem ``rank intrínseco'' baixo.
      \item \textbf{Por que $\mathbf{B}$ começa em zero?} Para que no início do treino $\Delta\mathbf{W} = \frac{\alpha}{r} \cdot \mathbf{0} \cdot \mathbf{A} = \mathbf{0}$: o modelo começa exatamente do pré-treinamento, sem perturbação. Se inicializássemos ambos aleatoriamente, a primeira forward pass divergiria do modelo base, perdendo o ponto de partida calibrado. Boa pergunta para a turma: por que não zerar $\mathbf{A}$ em vez de $\mathbf{B}$? (Se ambos fossem zero, o gradiente de ambos seria zero e nada seria aprendido; zerar só um basta.)
      \item \textbf{Contagem de parâmetros:} $r(d+k)$ vs.\ $dk$. Para uma projeção com $d=k=4096$, $r=16$: LoRA usa $16\times8192=131\,072$ params vs.\ $16.7$M do peso original (0,78\%). Para LLaMA-7B inteiro, Q e V em 32 camadas: $2 \times 32 \times 131\,072 \approx 8.4$M, ou $\sim$0,12\% de 6,7B.
      \item O fator $\alpha/r$ é análogo a uma taxa de aprendizado adicional para os adaptadores. Mantê-lo constante ao variar $r$ é uma boa prática (ex.: sempre use $\alpha = 2r$).
      \item \textbf{Deploy:} antes de servir, funde-se $\mathbf{W}_0 + \frac{\alpha}{r}\mathbf{B}\mathbf{A}$ num único tensor: sem custo extra de inferência em relação ao modelo base!
    \end{itemize}
  }
  \end{frame}

  %--- Slide: LoRA - Diagrama de Decomposição ---
  \begin{frame}{LoRA: Visualização da Decomposição \textit{Low Rank}}
    \begin{columns}[T]
      \begin{column}{0.58\textwidth}
        \centering
        \textcolor{DarkBlue}{\textbf{Atualização de pesos: modelo congelado + adaptadores}}
        \vspace{0.3em}

        \begin{tikzpicture}[every node/.style={font=\small}]
          \def\D{2.8}
          \def\K{2.4}
          \def\R{0.45}

          % W0 matrix (frozen)
          \filldraw[fill=LightGray, draw=MidGray, line width=1pt]
            (0, 0) rectangle (\K, \D);
          \node[font=\normalsize\bfseries] at ({\K/2}, {\D/2+0.15}) {$\mathbf{W}_0$};
          \node[font=\scriptsize, text=MidGray] at ({\K/2}, {\D/2-0.25})
            {\textit{congelado}};
          \draw[decorate, decoration={brace,amplitude=5pt}, DarkBlue]
            (0, \D+0.1) -- (\K, \D+0.1)
            node[midway, above=4pt, font=\scriptsize, text=DarkBlue] {$k$};
          \draw[decorate, decoration={brace,amplitude=5pt}, DarkBlue]
            (-0.1, \D) -- (-0.1, 0)
            node[midway, left=4pt, font=\scriptsize, text=DarkBlue] {$d$};

          % Plus
          \node[font=\LARGE\bfseries] at ({\K+0.55}, {\D/2}) {$+$};

          % B matrix (d x r, tall and thin)
          \pgfmathsetmacro{\bx}{\K+1.15}
          \filldraw[fill=CellBlue, draw=DarkBlue, line width=1pt]
            (\bx, 0) rectangle ({\bx+\R}, \D);
          \node[rotate=90, font=\small\bfseries, text=DarkBlue]
            at ({\bx+\R/2}, {\D/2}) {$\mathbf{B}$};
          \draw[decorate, decoration={brace,amplitude=4pt}, DarkBlue]
            (\bx, \D+0.1) -- ({\bx+\R}, \D+0.1)
            node[midway, above=4pt, font=\scriptsize, text=DarkRed] {$r \ll k$};

          % Times
          \node[font=\large] at ({\bx+\R+0.38}, {\D/2}) {$\times$};

          % A matrix (r x k, short and wide)
          \pgfmathsetmacro{\ax}{\bx+\R+0.80}
          \filldraw[fill=PastelGreen, draw=DarkGreen, line width=1pt]
            (\ax, {\D/2-\R/2}) rectangle ({\ax+\K}, {\D/2+\R/2});
          \node[font=\small\bfseries, text=DarkGreen]
            at ({\ax+\K/2}, {\D/2}) {$\mathbf{A}$};

          % Formula at bottom
          \node[font=\normalsize] at ({(\K + \bx+\R+0.80+\K)/2}, {-0.75})
            {$\mathbf{h} = \underbrace{\mathbf{W}_0\,\mathbf{x}}_{\text{congelado}} +
              \dfrac{\alpha}{r}\;\underbrace{\mathbf{B}\,\mathbf{A}\,\mathbf{x}}_{\text{treinável}}$};
        \end{tikzpicture}
      \end{column}
      \begin{column}{0.40\textwidth}
        \vspace{1.0em}
        \textcolor{DarkBlue}{\textbf{Contagem de parâmetros}}\\[0.3em]
        {\scriptsize modelo 7B: $d{=}k{=}4096$, $r{=}16$}
        \vspace{0.3em}

        \resizebox{\linewidth}{!}{%
        \begin{tabular}{lcc}
          \toprule
          \textbf{Matriz} & \textbf{Dim.} & \textbf{Params} \\
          \midrule
          \rowcolor{LightGray}
          $\mathbf{W}_0$ & $d \times k$ & 16.7M \textit{(frozen)} \\
          $\mathbf{B}$   & $d \times r$ & 65k \\
          $\mathbf{A}$   & $r \times k$ & 65k \\
          \midrule
          \multicolumn{2}{l}{\textbf{Total treinável}} & \textbf{130k} \\
          \multicolumn{2}{l}{\textcolor{DarkGreen}{\small fração}} &
            \textcolor{DarkGreen}{\textbf{$\approx$0,8\%}} \\
          \bottomrule
        \end{tabular}}

        \vspace{0.5em}
        \textcolor{DarkRed}{\scriptsize \textbf{Deploy:} fundir
        $\mathbf{W}_0 + \tfrac{\alpha}{r}\mathbf{B}\mathbf{A}$ antes de servir $\to$ zero overhead
        de inferência.}
      \end{column}
    \end{columns}
  \note{
    \textbf{Como ler o diagrama:}
    \begin{itemize}
      \item $\mathbf{W}_0$ (cinza) tem dimensões $d \times k$: por exemplo, uma projeção de atenção com $d=k=4096$ tem 16,7M parâmetros. Ele é congelado (gradiente bloqueado).
      \item $\mathbf{B}$ (azul, $d \times r$) e $\mathbf{A}$ (verde, $r \times k$) são os adaptadores de baixo rank. Com $r=16$: $\mathbf{B}$ tem $4096 \times 16 = 65536$ params e $\mathbf{A}$ tem $16 \times 4096 = 65536$ params (total de 131k, 0,78\% do peso original).
      \item $\mathbf{A}$ inicializa com distribuição normal gaussiana; $\mathbf{B}$ inicializa com zero: garante que $\Delta\mathbf{W} = \frac{\alpha}{r}\mathbf{B}\mathbf{A} = \mathbf{0}$ no início, preservando o modelo pré-treinado intacto.
      \item No deploy: funde-se $\mathbf{W}_0 + \frac{\alpha}{r}\mathbf{B}\mathbf{A}$ em um único tensor antes de servir (zero overhead em inferência).
    \end{itemize}
  }
  \end{frame}

  %--- Slide: QLoRA ---
  \begin{frame}{QLoRA: LoRA com Base Quantizada em 4 Bits \cite{dettmers2023qlora}}
    \begin{columns}
      \begin{column}{0.52\textwidth}
        Dettmers et al.\ (2023) combinam LoRA com quantização do modelo base,
        permitindo ajustar modelos de 65B em uma única GPU de 48~GB.

        \begin{block}{Componentes chave do QLoRA}
          \begin{enumerate}
            \item \textbf{NF4} (\textit{NormalFloat 4-bit}): minimiza o erro
              de quantização para pesos com distribuição normal (vs.\ INT4 uniforme)
            \item \textbf{Dupla quantização}: quantizar também as constantes
              de quantização, economizando $\approx$0,37 bits/parâmetro extra
            \item \textbf{\textit{Paged Optimizers}}: memória unificada CPU/GPU absorve picos de uso
          \end{enumerate}
        \end{block}
      \end{column}
      \begin{column}{0.45\textwidth}
        \textcolor{DarkBlue}{\textbf{Comparação de memória (modelo 7B)}}
        \vspace{0.3em}

        \resizebox{\linewidth}{!}{%
        \begin{tabular}{lcc}
          \toprule
          \textbf{Método}     & \textbf{VRAM (treino)} & \textbf{Precisão} \\
          \midrule
          Full FT (fp16)      & $\sim$112 GB & fp16 \\
          LoRA (fp16)         & $\sim$40 GB  & fp16 \\
          QLoRA (NF4)         & $\sim$10 GB  & NF4 + fp16 \\
          \bottomrule
        \end{tabular}}

        \vspace{0.4em}
        \textcolor{DarkGreen}{\small \textbf{Custo:} pequena degradação de
        qualidade ($\approx$0,1--0,2 pontos em benchmarks) em troca de
        acessibilidade em hardware de consumo.}

        \vspace{0.3em}
        \textcolor{MidGray}{\small \textbf{Implementação:} bibliotecas
        \texttt{bitsandbytes} + \texttt{peft} do HuggingFace.}
      \end{column}
    \end{columns}
  \note{
    \textbf{Intuição do QLoRA e seus três componentes:}
    \begin{itemize}
      \item \textbf{NF4 (NormalFloat 4-bit):} pesos de redes neurais seguem distribuição aproximadamente normal. O NF4 usa 16 valores de quantização distribuídos de acordo com a \emph{função de distribuição acumulada} de $\mathcal{N}(0,1)$: os bins são mais densos perto do zero (onde há mais pesos) e mais esparsos nas caudas. Isso minimiza o erro quadrático médio de quantização comparado ao INT4 uniforme.
      \item \textbf{Dupla quantização:} normalmente, cada grupo de 64 pesos tem uma constante de escala em fp32 (32 bits para 64 params = 0,5 bits extras/param). A dupla quantização quantiza essas constantes de escala em fp8 (8 bits), economizando $\approx 0,37$ bits/param adicionais.
      \item \textbf{Paged Optimizers:} durante gradient checkpointing e backward pass, há picos de uso de memória. O NVIDIA Unified Memory permite ``paginar'' o estado do otimizador para RAM da CPU quando a GPU satura, evitando OOM (Out of Memory) errors.
      \item \textbf{Ponto crucial:} apenas os pesos base ($\mathbf{W}_0$) são quantizados em NF4 e \emph{congelados}. Os adaptadores $\mathbf{A}$ e $\mathbf{B}$ ficam em bf16 e são treináveis normalmente. Os gradientes fluem pelos adaptadores sem precisar desquantizar $\mathbf{W}_0$.
    \end{itemize}
  }
  \end{frame}

  %--- Slide: NF4 e Quantização para PEFT ---
  \begin{frame}{Quantização para PEFT: NF4 vs.\ INT4}
    \begin{block}{NormalFloat 4-bit (NF4)}
      \small
      Tipo de dado de 4 bits projetado para pesos de redes neurais, que
      tipicamente seguem $\mathcal{N}(0, \sigma^2)$. Em vez de espaçamento
      uniforme (INT4), os 16 valores representáveis do NF4 são obtidos dos
      \textbf{quantis da distribuição normal}: cada nível é a média de dois
      quantis vizinhos de $\mathcal{N}(0,1)$ (o que evita os extremos $\pm\infty$),
      com um zero exato, e o conjunto é normalizado para $[-1, 1]$. Mais
      resolução perto de zero, onde há mais pesos; menos resolução nas caudas.
    \end{block}

    \vspace{0.3em}
    \centering
    \begin{tikzpicture}[yscale=1.0, xscale=1.9]
      % Normal density (scaled to peak ~1.5)
      \fill[DarkBlue!15, domain=-2.7:2.7, samples=80]
        plot (\x, {1.5*exp(-\x*\x/2)}) -- (2.7, 0) -- (-2.7, 0) -- cycle;
      \draw[DarkBlue, thick, domain=-2.7:2.7, samples=80]
        plot (\x, {1.5*exp(-\x*\x/2)});
      \node[font=\scriptsize, text=DarkBlue, anchor=west] at (1.0, 1.35)
        {$\mathcal{N}(0,1)$: distribuição dos pesos};

      % Axis
      \draw[->, thick] (-2.8, 0) -- (2.9, 0)
        node[right, font=\tiny] {$w$};
      \foreach \x/\lbl in {-2/$-2\sigma$, -1/$-\sigma$, 0/$0$, 1/$+\sigma$, 2/$+2\sigma$} {
        \draw[thick] (\x, -0.05) -- (\x, 0.05);
        \node[font=\tiny, below=1pt] at (\x, -0.05) {\lbl};
      }

      % NF4 quantile ticks (above axis, green)
      \foreach \x in {-2.40,-1.67,-1.26,-0.95,-0.68,-0.44,-0.22,0.00,0.19,0.39,0.59,0.81,1.06,1.35,1.74,2.40} {
        \draw[DarkGreen, very thick] (\x, 0.02) -- (\x, 0.32);
      }
      \node[font=\scriptsize\bfseries, text=DarkGreen, anchor=west]
        at (-2.8, 0.95) {NF4: \textit{denso no centro}};

      % INT4 uniform ticks (below axis, red)
      \foreach \x in {-2.40,-2.08,-1.76,-1.44,-1.12,-0.80,-0.48,-0.16,0.16,0.48,0.80,1.12,1.44,1.76,2.08,2.40} {
        \draw[DarkRed, very thick] (\x, -0.32) -- (\x, -0.02);
      }
      \node[font=\scriptsize\bfseries, text=DarkRed, anchor=west]
        at (-2.6, -0.55) {INT4: \textit{uniforme, desperdiça nas caudas}};
    \end{tikzpicture}

    \vspace{0.4em}
    \textcolor{DarkGreen}{\small \textbf{QLoRA:} pesos base em NF4 (congelados);
    adaptadores LoRA em bf16 (treináveis). Menor erro de quantização que INT4
    com os mesmos 4 bits.}
  \note{
    \textbf{Intuição do NF4 vs INT4:}
    \begin{itemize}
      \item INT4 uniforme divide $[\min(W), \max(W)]$ em 16 bins iguais. Como a distribuição de pesos é normal (concentrada perto do zero), a maioria dos valores cai em apenas 2--4 bins centrais, desperdiçando resolução nas caudas.
      \item NF4 usa os quantis de $\mathcal{N}(0,1)$: coloca mais bins onde há mais massa de probabilidade (perto do zero) e menos nas caudas raras, minimizando o erro quadrático médio de quantização. \textbf{Construção exata (Dettmers et al., 2023):} para $k=4$ bits, tomam-se $2^{k-1}$ quantis para a metade negativa e $2^{k-1}+1$ para a positiva, cada nível sendo a média de dois quantis consecutivos (assim os extremos ficam finitos), unem-se as duas metades removendo um dos zeros duplicados, e normaliza-se para $[-1,1]$. Cada bloco de 64 pesos é dividido pelo seu máximo absoluto antes de ser mapeado nesses 16 níveis.
      \item Na figura, as marcas verdes são os 16 níveis exatos do NF4 (os valores da tabela do \texttt{bitsandbytes}, em $[-1,1]$) multiplicados por $2{,}4\sigma$, que é o \textit{absmax} esperado de um bloco de 64 pesos gaussianos; as vermelhas são 16 níveis uniformes no mesmo intervalo. Note a assimetria do NF4: 7 níveis negativos, o zero e 8 positivos, porque o zero exato ``rouba'' um nível de um dos lados. O espaçamento vai de $\approx 0{,}08$ perto do zero a $\approx 0{,}3$ nas caudas (em unidades normalizadas).
      \item \textbf{Ponto crucial:} apenas os pesos base ($\mathbf{W}_0$) são quantizados em NF4 e congelados. Os adaptadores $\mathbf{A}$ e $\mathbf{B}$ ficam em bf16: os gradientes fluem pelos adaptadores sem precisar desquantizar $\mathbf{W}_0$ para o otimizador (a desquantização acontece só no forward, bloco a bloco).
    \end{itemize}
  }
  \end{frame}

  %--- Slide: Tabela Comparativa PEFT ---
  \begin{frame}{Comparação dos Métodos PEFT}
    \centering
    \resizebox{0.98\linewidth}{!}{%
    \begin{tabular}{lccccp{3.5cm}}
      \toprule
      \textbf{Método} &
      \textbf{Params. treináveis} &
      \textbf{VRAM (7B)} &
      \textbf{Desempenho} &
      \textbf{Complexidade} &
      \textbf{Caso de uso típico} \\
      \midrule
      \rowcolor{CellYellow}
      Full FT       & 100\%      & $\sim$112 GB & Máximo     & Baixa  &
        Dados abundantes, hardware dedicado \\
      LoRA          & 0,1--1\%   & $\sim$40 GB  & Alto       & Baixa  &
        Ajuste geral, produção \\
      \rowcolor{CellBlue}
      QLoRA         & 0,1--1\%   & $\sim$10 GB  & Alto$^*$   & Média  &
        Hardware de consumo (GPU única) \\
      \bottomrule
      \multicolumn{6}{l}{\small $^*$ Pequena degradação vs.\ LoRA em fp16}
    \end{tabular}}

    \vspace{0.4em}
    \begin{block}{Recomendação prática}
      Para a maioria dos cenários, \textbf{QLoRA} é o padrão para ajuste em
      hardware limitado; \textbf{LoRA em bf16} quando há mais VRAM disponível.
      \textit{Full fine-tuning} apenas quando há hardware dedicado e dados abundantes.
    \end{block}
  \note{
    \textbf{Como navegar esta tabela comparativa:}
    \begin{itemize}
      \item A linha ``Full FT'' é a baseline: máximo desempenho, máxima memória. Tudo abaixo é um trade-off.
      \item \textbf{LoRA $\sim$40 GB:} o modelo base está em fp16 (14 GB), mas ainda precisa de gradientes para backward pass através das camadas congeladas (para chegar aos adaptadores): isso consome memória. Técnicas como gradient checkpointing ajudam.
      \item \textbf{QLoRA $\sim$10 GB:} o modelo base está em NF4 ($\sim$3,5 GB para 7B), os adaptadores + gradientes em bf16 ($\sim$6,5 GB). Cabe numa RTX 3060 12GB ou T4 do Colab.
      \item Recomendação de ordem de escolha: QLoRA $\to$ LoRA $\to$ Full FT, à medida que o hardware disponível aumenta.
    \end{itemize}
  }
  \end{frame}

  %======================================================================================
  % LEITURAS RECOMENDADAS
  %======================================================================================
  \section*{Referências e Leituras}

  \begin{frame}{Leituras Recomendadas}
    \leiturasize
    \begin{itemize}
      \item Lambert, N. (2025). \textit{Reinforcement Learning from Human Feedback} (The RLHF Book). \textit{rlhfbook.com}. RM, PPO, DPO, GRPO e RLVR em um só lugar.
      \item Hugging Face (2025). \textit{The Ultra-Scale Playbook: Training LLMs on GPU Clusters}. Paralelismo, ZeRO e precisão mista com experimentos.
      \item Grattafiori, A. et al. (2024). The Llama 3 Herd of Models. \textit{arXiv:2407.21783}. Relato completo de pré-treino e pós-treino.
      \item Lambert, N. et al. (2024). Tülu 3: Pushing Frontiers in Open Language Model Post-Training. \textit{arXiv:2411.15124}.
      \item Shao, Z. et al. (2024). DeepSeekMath. \textit{arXiv:2402.03300}. Origem do GRPO.
      \item Penedo, G. et al. (2024). The FineWeb Datasets: Decanting the Web for the Finest Text Data at Scale. \textit{NeurIPS Datasets and Benchmarks}.
      \item Karpathy, A. (2024). Let's Reproduce GPT-2 (124M). \textit{github.com/karpathy/build-nanogpt}. Pré-treino completo em código.
    \end{itemize}
  \end{frame}

  %======================================================================================
  % REFERÊNCIAS
  %======================================================================================
  \begin{frame}[allowframebreaks]{Referências}
    \sloppy
    \printbibliography[heading=none]
  \end{frame}

\end{document}
